\documentclass[12pt,reqno]{amsart}
\usepackage{mathrsfs}
\usepackage{cases}
\usepackage{epic}
\usepackage{amsfonts}
\usepackage{graphicx}
\usepackage{amsmath}
\usepackage{amssymb, upgreek}
\usepackage{bm}
\usepackage{latexsym,todonotes}
\usepackage{pdflscape}
\usepackage[all]{xy}
\usepackage{color}
\usepackage{colordvi}
\usepackage{multicol}
\usepackage{enumitem}
\usepackage[normalem]{ulem}
\textwidth=14.5cm \oddsidemargin=1cm \evensidemargin=1cm
\setlength{\hoffset}{0pt}
\setlength{\voffset}{0pt}
\setlength{\topmargin}{0pt}
\setlength{\oddsidemargin}{0in}
\setlength{\evensidemargin}{0in}
\setlength{\textheight}{8.75in}
\setlength{\textwidth}{6.5in}
\pagestyle{headings}
\newcommand{\arxiv}[1]{\href{http://arxiv.org/abs/#1}{\tt arXiv:\nolinkurl{#1}}}


\usepackage{tikz-cd}
\numberwithin{equation}{section}
\usepackage[linktocpage=true]{hyperref}
\hypersetup{colorlinks,linkcolor=blue,urlcolor=cyan,citecolor=blue}
\usetikzlibrary{arrows.meta}


\newtheorem{theorem}{Theorem}[section]
\newtheorem{proposition}[theorem]{Proposition}
\newtheorem{lemma}[theorem]{Lemma}
\newtheorem{corollary}[theorem]{Corollary}
\newtheorem{conjecture}[theorem]{Conjecture}
\theoremstyle{definition}
\newtheorem{definition}[theorem]{Definition}
\newtheorem{example}[theorem]{Example}
\theoremstyle{remark}
\newtheorem{remark}[theorem]{Remark}

\newcommand{\C}{\mathbb C}
\newcommand{\N}{\mathbb N}
\newcommand{\Z}{\mathbb Z}
\newcommand{\A}{\mathscr A}
\newcommand{\M}{\mathcal M}
\newcommand{\U}{\mathbf U}
\newcommand{\nplus}{\mathfrak n_+}
\newcommand{\g}{\mathfrak g}
\newcommand{\Flag}{\operatorname{Flag}}
\newcommand{\rk}{\rank}
\newcommand{\ad}{\operatorname{ad}}
\newcommand{\wt}{\operatorname{wt}}
\newcommand{\htop}{\operatorname{ht}}
\newcommand{\Hom}{\operatorname{Hom}}
\newcommand{\Ext}{\operatorname{Ext}}
\newcommand{\End}{\operatorname{End}}
\newcommand{\Spec}{\operatorname{Spec}}
\newcommand{\Sym}{\operatorname{Sym}}
\newcommand{\im}{\operatorname{im}}
\newcommand{\diag}{\operatorname{diag}}
\newcommand{\lcm}{\operatorname{lcm}}
\newcommand{\eps}{\varepsilon}
\newcommand{\un}{\mathbf 1}
\newcommand{\stmod}{\underline{\operatorname{mod}}}
\newcommand{\sgn}{\operatorname{sgn}\nolimits}
\renewcommand{\Im}{\operatorname{Im}}
\newcommand{\Ker}{\operatorname{Ker}}
\newcommand{\rank}{\operatorname{rank}}
\newcommand{\Rep}{\operatorname{Rep}}
\newcommand{\Gr}{\operatorname{Gr}}
\newcommand{\Fl}{\operatorname{Fl}}
\newcommand{\GL}{\operatorname{GL}}


\theoremstyle{plain}
\newtheorem{introtheorem}{Theorem}
\renewcommand{\theintrotheorem}{\Alph{introtheorem}}
\theoremstyle{plain}
\newtheorem*{maintheorem}{Main theorem}

\title[Generalized Preprojective Convolution Algebras]
{Convolution Algebras of Generalized Preprojective Algebras and Kac--Moody Positive Parts}
\author{Fengjun Shi}
\address{Fengjun Shi, Department of Mathematics, Sichuan University, Chengdu 610064, P.~R.~China}
\email{shifengjun@stu.scu.edu.cn}
\thanks{This work was partially supported by the National Natural
Science Foundation of China (Grant No.~12471037).}
\date{}
\keywords{Generalized preprojective algebra, constructible function, convolution algebra, Kac--Moody algebra}



\begin{document}

\begin{abstract}
We realize the enveloping algebras of positive parts of symmetrizable Kac--Moody Lie algebras as convolution algebras of constructible functions on module varieties of Geiss--Leclerc--Schr\"oer generalized preprojective algebras $\Pi(C,D)$.
% We define a symmetrizable generalized Cartan matrix $A=A(C,D)$ and characterize when the canonical homomorphism $U(\mathfrak n_+(A))\to\mathcal M(C,D)$ is an isomorphism. 
We characterize all pairs $(C,D)$ for which such a realization exists.
These realizations generalize Lusztig's classical preprojective construction.
The construction uses a folding embedding and torus actions on flag varieties.
\end{abstract}

    \maketitle
    \setcounter{tocdepth}{1}
    \tableofcontents

%%%%%%%%%%%%%%%%%%%%%%%%%%%%%%%%%%%%%%%%%%%%%%
\section{Introduction}\label{sec:introduction}
%\subsection{Hall algebras and geometric realizations}
Hall algebras provide representation-theoretic constructions of Lie algebras and their enveloping algebras. For a symmetrizable Kac--Moody Lie algebra $\mathfrak g$, write $\mathbf{U}_v^+(\g)$ for the positive part of the corresponding quantum group. Using representations of hereditary algebras over finite fields, Ringel established a Hall algebra realization of $\mathbf U_v^+(\g)$ in finite type, and Green extended it to the general symmetrizable case \cite{Rin90,Gr95}. The classical specialization at $v=1$ gives $U(\nplus)$, where $\nplus$ is the positive part of $\mathfrak g$.

In the symmetric case, Schofield \cite{Sch91} realized $U(\nplus)$ over $\C$ by replacing finite-field point counts with Euler characteristics of complex varieties. Lusztig \cite{Lus91} reformulated this construction using convolution of constructible functions on quiver representation varieties. He also realized $U(\nplus)$ by convolution on varieties of nilpotent representations of classical preprojective algebras \cite{Lus91,Lus00}.

%Riedtmann \cite{Rie94} used Euler-characteristic convolution to construct Lie algebras and realize their enveloping algebras from finite-dimensional complex algebras of finite representation type. Ding, Xiao, and Xu \cite{DXX10} extended this approach to arbitrary finite-dimensional complex algebras. 

%For quivers with loops, Bozec \cite{Boz15} realized the positive parts of the associated quantum Borcherds-Bozec algebras using perverse sheaves. For the classical specialization at $q=1$, he obtained a constructible-function realization and constructed semicanonical bases \cite{Boz16}. Hennecart \cite{Hen24} compared and generalized these geometric realizations for arbitrary quivers and established their integral versions. In particular, he obtained a realization using the top cohomological Hall algebra of the moduli stack of strictly seminilpotent representations of the associated preprojective algebra, with a suitable sign twist of the multiplication.

%The representation theory of classical preprojective algebras also plays an important role in the study of coordinate rings and cluster algebras. Geiss, Leclerc, and Schr\"oer studied the multiplicative properties of the dual semicanonical basis \cite{GLS05} and established multiplication formulas for evaluation forms in terms of extension varieties \cite{GLS07}. In simply laced Dynkin type, they used maximal rigid modules and their mutations to categorify the cluster algebra structure on $\mathbb C[N]$, where $N$ is a maximal unipotent subgroup of the corresponding complex semisimple group, and proved that all cluster monomials belong to the dual semicanonical basis \cite{GLS06}. They subsequently extended this approach to the coordinate rings of unipotent subgroups and unipotent cells of symmetric Kac-Moody groups \cite{GLS11}.

%\subsection{Generalized preprojective algebras}
To extend these geometric constructions to the symmetrizable setting, Geiss, Leclerc, and Schr\"oer \cite{GLS17} introduced the algebras $H(C,D,\Omega)$ and the generalized preprojective algebras $\Pi(C,D)$, where $C$ is a symmetrizable generalized Cartan matrix, $D$ is a positive integral symmetrizer, and $\Omega$ is an acyclic orientation of $C$. Using locally free $H(C,D,\Omega)$-modules, they obtained a convolution realization of $U(\mathfrak n_+(C))$ when $C$ is of Dynkin type \cite{GLS16}. Castillo G\'omez and Geiss \cite{CG25} proved the corresponding result in affine type $\widetilde C_n$ with the minimal symmetrizer. 

%For $\Pi(C,D)$, they extended the geometric construction of crystal graphs to the symmetrizable setting and constructed semicanonical functions \cite{GLS18}.

For $\Pi(C,D)$, let $\mathcal M(C,D)$ denote the unital convolution subalgebra generated by the characteristic functions $\theta_i=\mathbf{1}_{E_i}$ of the generalized simple $\Pi(C,D)$-modules $E_i$. These generators need not satisfy the Serre relations associated with $C$. In \cite{GLS18}, the authors considered the quotient by the ideal generated by the corresponding Serre elements and proved that it is isomorphic to $U(\nplus(C))$ under their Support Conjecture.


Fei \cite[Theorem~10.1]{Fei26} proved this isomorphism without assuming the Support Conjecture. His proof uses the convolution realization for $H(C,D,\Omega)$, which he established for all symmetrizable $C$ and all positive integral symmetrizers $D$ \cite[Theorem~A.1]{Fei26}. He also gave a counterexample to the conjecture \cite[Theorem~10.2]{Fei26}. In the same work, he realized $U(\nplus(C))$ using monodromy traces on flag varieties for a ramified preprojective species over $\C((t))$.

\subsection{Main results}
We determine when the convolution algebra $\M(C,D)$ is isomorphic to the enveloping algebra of the positive part of a symmetrizable Kac--Moody Lie algebra.

Write $C=(c_{ij})_{i,j\in I}$ and $D=\diag(d_i)_{i\in I}$, and set
\[
    \lambda_i=
    \begin{cases}
        1, & d_i=1,\\
        2, & d_i>1,
    \end{cases}
\]
and define $A=A(C,D)=(a_{ij})_{i,j\in I}$ by
\[
    a_{ii}=2,\qquad
    a_{ij}=\lambda_i c_{ij}\quad(i\ne j).
\]
Then $A$ is a symmetrizable generalized Cartan matrix, with positive integral symmetrizer $\diag(2d_i/\lambda_i)_{i\in I}$. 
Let $e_i$ be the Chevalley generators of $\nplus (A)$, and let $\mathfrak l(C,D)$ be the Lie subalgebra of $\M(C,D)$ generated by $\theta_i$ under the convolution commutator.

For arbitrary $(C,D)$, we prove the $A$-Serre relations in Proposition~\ref{prop:A-Serre}. 
They give a canonical surjective Hopf algebra homomorphism
\[
    \Phi_{C,D}\colon U(\nplus (A)) \twoheadrightarrow \M(C,D), \qquad e_i\mapsto\theta_i.
\]

We consider the following condition, denoted by condition~\eqref{eq:edge-condition}:
\[
    \gcd(d_i,d_j)>1
    \qquad\text{whenever }c_{ij}<0
    \text{ and }d_i,d_j>1.
\]

\begin{introtheorem}[Theorem~\ref{thm:main}]
\label{thm:introduction-main}
If $(C,D)$ satisfies condition~\eqref{eq:edge-condition}, then  
$$\Phi_{C,D}:
    U(\nplus (A))
    \rightarrow
    \M(C,D)$$ 
    is an isomorphism of rank-graded Hopf algebras. Its restriction gives a Lie algebra isomorphism
\[
    \nplus(A)
    \xrightarrow{\;\sim\;}
    \mathfrak l(C,D).
\]
\end{introtheorem}

When $D$ is the identity matrix, we have $A=C$, and Theorem~\ref{thm:introduction-main} recovers Lusztig's classical preprojective realization; see \cite[Section~1.2]{GLS18}. The theorem also applies whenever $C$ is of finite or affine type and $D$ is any positive integral symmetrizer, by Corollary~\ref{cor:finite-affine}. The matrix $A(C,D)$ need not have the same type as $C$.

For pairs $(C,D)$ not satisfying condition~\eqref{eq:edge-condition}, we prove the following result in Section~\ref{app:serre-quotients}.

\begin{introtheorem}[Theorem~\ref{thm:nonrealization}]
\label{thm:introduction-nonrealization}
If $(C,D)$ does not satisfy condition~\eqref{eq:edge-condition}, then
\[
    \mathcal M(C,D)\not\cong U(\mathfrak n_+(B))
\]
as unital $\C$-algebras for any symmetrizable generalized Cartan matrix $B$.
In particular, $\mathfrak l(C,D)$ is not isomorphic to the positive part of any symmetrizable Kac--Moody Lie algebra.
\end{introtheorem}

To obtain realizations for arbitrary $(C,D)$, we construct Serre quotients of $\M(C,D)$ in Section~\ref{app:serre-quotients}. In particular, we recover the full $C$-Serre quotient realization proved by Fei \cite[Theorem~10.1]{Fei26} without assuming the GLS Support Conjecture.

\subsection{Organization}
The paper is organized as follows. 
Section~\ref{sec:preliminaries} recalls generalized preprojective algebras and their convolution algebras, defines the Cartan matrix $A(C,D)$, and states the main theorems.
In Section~\ref{sec:surjectivity}, we prove a splitting lemma for locally free modules and use it to establish the $A$-Serre relations. This gives the canonical homomorphism $U(\mathfrak n_+(A))\to\mathcal M(C,D)$ for arbitrary $(C,D)$.
In Section~\ref{sec:injectivity}, we use a symmetric covering of $A$ and torus actions on flag varieties to prove injectivity under condition ~\eqref{eq:edge-condition}; the flag comparison is related to the acyclic construction in \cite[Appendix~A]{Fei26}.
Section~\ref{app:serre-quotients} proves Theorem~\ref{thm:nonrealization} and constructs Serre quotients for arbitrary $(C,D)$.

\medskip
\noindent\textbf{Acknowledgements.}
The author would like to thank Changjian Fu, Shengfei Geng and Ming Lu for helpful discussions. 
%%%%%%%%%%%%%%%%%%%%%%%%%%%%%%%%%%%%%%%%%%%%%%%
\section{Generalized preprojective algebras and convolution algebras}
\label{sec:preliminaries}
We work over $\C$, and all modules are finite-dimensional left modules.
\subsection{Symmetrizable Cartan matrices and orientations}
    We recall the definitions from \cite{GLS17}. Set $I=\{1,\ldots,n\}$ and $\N=\{0,1,2,\ldots\}$.

    %Following \cite{GLS17}, we recall basic definitions and properties of the preprojective algebras associated with symmetrizable generalized Cartan matrices. Set $I=\{1,\ldots,n\}$ and $\N=\{0,1,2,\ldots\}$.

    \begin{definition}A matrix $C = (c_{ij}) \in M_n(\mathbb{Z})$ is a symmetrizable (generalized) Cartan matrix if the following conditions are satisfied:
     \begin{itemize}
         \item[($C_1$)] $c_{ii} = 2$ for all $i$;
         \item[($C_2$)] $c_{ij} \leq 0$ for all $i \neq j$;
         \item[($C_3$)] $c_{ij} \neq 0$ if and only if $c_{ji} \neq 0$;
         \item[($C_4$)] There is a diagonal integer matrix $D = \diag(d_1,\cdots,d_n)$ with $d_i \geq 1$ for all $i$ such that $DC$ is symmetric.
     \end{itemize}
     \end{definition}
      The matrix $D$ in $(C_4)$ is called a symmetrizer of $C$. It is minimal if $d_1 +\cdots +d_n$ is minimal. For each pair with $c_{ij}<0$, set
     \begin{align*}
         g_{ij}:=|\gcd(c_{ij},c_{ji})|, \quad f_{ij}:=|c_{ij}|/g_{ij}, \quad k_{ij}:= \gcd(d_{i},d_{j}).
     \end{align*}
     Since $d_if_{ij}=d_jf_{ji}$ and $\gcd(f_{ij},f_{ji})=1$, we have
     \[
     d_i=k_{ij}f_{ji}, \qquad d_j=k_{ij}f_{ij}.
     \]

     An orientation of $C$ is a subset $\Omega \subseteq I\times I$ satisfying the following conditions:
     \begin{enumerate}[label=(\roman*)]
         \item $\{(i, j), (j, i)\} \cap \Omega \neq \emptyset $ if and only if $c_{ij} < 0$;
         \item For each sequence $((i_1, i_2), (i_2, i_3), \cdots , (i_t, i_{t+1}))$ with $t \geq 1$ and $(i_s, i_{s+1}) \in \Omega$ for all $1 \leq s \leq t$, we have $i_1 \neq i_{t+1}$.
     \end{enumerate}

     For an orientation $\Omega$ of $C$, let $Q=Q(C,\Omega)$ be the quiver with vertex set $Q_0=I$ and arrow set
     \[
         Q_1= 
    \left\{
    \alpha_{ij}^{(g)}:j\to i \;\middle|\; (i,j)\in\Omega,\ 1\leq g\leq g_{ij} \right\}
    \cup 
    \left\{
    \varepsilon_i:i\to i \;\middle|\; i\in I \right\}.
     \]
     Thus, $Q$ has $g_{ij}$ arrows from $j$ to $i$ for each $(i,j)\in\Omega$, and one loop $\varepsilon_i$ at each vertex. The quiver obtained by deleting all loops $\varepsilon_i$ is acyclic.

     Set
     \[
         \Omega^{\mathrm{op}} =\{(j,i)\mid(i,j)\in\Omega\}, \qquad \overline{\Omega}=\Omega\cup\Omega^{\mathrm{op}}.
     \]

    Let $\overline{Q}$ be the quiver obtained from $Q$ by adding a reverse arrow $\alpha_{ji}^{(g)}:i\to j$ for each non-loop arrow $\alpha_{ij}^{(g)}:j\to i$. The loops $\varepsilon_i$ are retained without adding new copies. For $(i,j)\in\overline{\Omega}$, define
    \[
    s_{ij}= \begin{cases}
    1, & (i,j)\in\Omega,\\
    -1, & (i,j)\in\Omega^{\mathrm{op}}.
    \end{cases}
    \]

\subsection{Generalized preprojective algebras}
    %We first recall the algebra $H(C,D,\Omega)$ associated with the oriented quiver $Q$, and then define the generalized preprojective algebra using $\overline{Q}$.

    \begin{definition}\label{def:H}
    The algebra $H=H(C,D,\Omega)$ is the quotient
    \[
    H(C,D,\Omega)=\mathbb{C}Q/\mathcal{I}_H,
    \]
    where $\mathcal{I}_H$ is the two-sided ideal generated by the following relations:
    \begin{enumerate}
    \item For each $i\in I$, the nilpotency relation
    \[
    \varepsilon_i^{d_i}=0.
    \]
    \item For each $(i,j)\in\Omega$ and
    $1\leq g\leq g_{ij}$, the commutativity relation
    \[
    \varepsilon_i^{f_{ji}}\alpha_{ij}^{(g)}
    =
    \alpha_{ij}^{(g)}\varepsilon_j^{f_{ij}}.
    \]
    \end{enumerate}
    \end{definition}

    %The generalized preprojective algebra is defined on $\overline{Q}$ by imposing these relations in both directions, together with a mesh relation at each vertex.

    \begin{definition}\label{def:Pi}
    The generalized preprojective algebra $\Pi=\Pi(C,D,\Omega)$ is the quotient
    \[
    \Pi(C,D,\Omega)
    =
    \mathbb{C}\overline{Q}/\mathcal{I}_{\Pi},
   \]
   where $\mathcal{I}_{\Pi}$ is the two-sided ideal generated by the following relations:
    \begin{enumerate}
    \item For each $i\in I$, the nilpotency relation
    \[
    \varepsilon_i^{d_i}=0.
    \]
    \item For each $(i,j)\in\overline{\Omega}$ and
    $1\leq g\leq g_{ij}$, the commutativity relation
    \[
    \varepsilon_i^{f_{ji}}\alpha_{ij}^{(g)}
    =
    \alpha_{ij}^{(g)}\varepsilon_j^{f_{ij}}.
    \]
    \item For each $i\in I$, the mesh relation
    \[
    \sum_{\substack{j\in I\\c_{ij}<0}}
    s_{ij}
    \sum_{g=1}^{g_{ij}}
    \sum_{p=0}^{f_{ji}-1}
    \varepsilon_i^p
    \alpha_{ij}^{(g)}\alpha_{ji}^{(g)}
    \varepsilon_i^{f_{ji}-1-p}
    =0.
    \]
    \end{enumerate}
    \end{definition}

    The isomorphism class of $\Pi(C,D,\Omega)$ is independent of $\Omega$, so we write $\Pi(C,D)$. If $C$ is symmetric and $D=I$, the loops vanish and $\Pi(C,D)$ is the ordinary preprojective algebra.

\subsection{The quadratic form}
    We associate with $(C,D)$ the quadratic form $q_{C,D}:\mathbb{Z}^n\to\mathbb{Z}$ defined by
    \[
    q_{C,D}(\mathbf{v})
    =
    \frac{1}{2}\mathbf{v}^{T}DC\mathbf{v}
    =
    \sum_{i=1}^{n}d_i v_i^2 - \sum_{i<j}d_i|c_{ij}|v_i v_j,
    \qquad
    \mathbf{v}=(v_1,\ldots,v_n).
    \]
    The matrix $C$ is of finite (Dynkin) type if and only if the extension of $q_{C,D}$ to $\mathbb{R}^n$ is positive definite. The connected finite types are
    \[
    A_n,\ B_n,\ C_n,\ D_n,\ E_6,\ E_7,\ E_8,\ F_4,\ G_2.
    \]
    If $C$ is connected and of affine type, this form is positive semidefinite but not positive definite.
    
\subsection{Locally free and generalized simple modules}
   Put $H_i=\C[\eps_i]/(\eps_i^{d_i})$. For a $\Pi(C,D)$-module $M$, write $M_i$ for its vector space at vertex $i$, regarded as an $H_i$-module via the action of $\varepsilon_i$. A module $M$ is \emph{locally free} if $M_i$ is free over $H_i$ for every $i$, and its rank vector is
   \[
   \rk M=(\rk_{H_i}M_i)_{i\in I}.
   \]
   The generalized simple module $E_i$ is supported at $i$ with component $H_i$. It is indecomposable; it is an ordinary simple module only when $d_i=1$.

   A module $M$ is called \emph{$E$-filtered} if it admits
   a finite filtration
   \[
   0=M^0\subset M^1\subset\cdots\subset M^m=M
   \]
   such that
   \[
   M^a/M^{a-1}\cong E_{i_a}, \qquad 1\leq a\leq m,
   \]
   for some $i_1,\ldots,i_m\in I$.

   Every $E$-filtered module is locally free. Indeed, at each vertex $j\in I$, the filtration induces short exact sequences
   \[
   0\longrightarrow (M^{a-1})_j \longrightarrow (M^a)_j \longrightarrow (E_{i_a})_j \longrightarrow 0.
   \]
   The quotient $(E_{i_a})_j$ is either $H_j$ or zero, so each sequence splits as a sequence of $H_j$-modules. It follows inductively that $M_j$ is free over $H_j$.

\subsection{The convolution algebra}
   Following  \cite[Section~4.1] {GLS16} and \cite[Section~2.5] {GLS18}, we recall the construction of the convolution algebra of constructible functions associated with $\Pi:=\Pi(C,D)$. 

   For an arrow $a$ of $\overline Q$, let $s(a)$ and $t(a)$ denote its source and target. For a dimension vector $\mathbf v=(v_i)_{i\in I}\in\mathbb N^I$, let $\Rep_{\mathbf v}(\Pi)$ be the affine variety of $\Pi$-module structures on the vertex spaces $(\mathbb C^{v_i})_{i\in I}$. Concretely, its points are tuples
   \[
   M=(M(a))_{a\in\overline Q_1} \in \prod_{a\in\overline Q_1} \Hom_{\mathbb C} \bigl(\mathbb C^{v_{s(a)}}, \mathbb C^{v_{t(a)}}\bigr)
   \]
   satisfying the defining relations of $\Pi$.

   The group
   \[
   G_{\mathbf v}
   =
   \prod_{i\in I}\GL_{v_i}(\mathbb C)
   \]
   acts on $\Rep_{\mathbf v}(\Pi)$ by changes of bases:
   \[
   (g\cdot M)(a)
   =
   g_{t(a)}M(a)g_{s(a)}^{-1}, \qquad g=(g_i)_{i\in I}\in G_{\mathbf v}.
   \]
   Its orbits correspond bijectively to isomorphism classes of $\Pi$-modules with dimension vector $\mathbf v$.

   A constructible function on a complex algebraic variety is a complex-valued function taking finitely many values and having constructible fibers. Let $\mathcal F_{\mathbf v}(\Pi)$ be the $\mathbb C$-vector space of constructible functions
   \[
   f:\Rep_{\mathbf v}(\Pi)\longrightarrow \mathbb C
   \]
   that are constant on $G_{\mathbf v}$-orbits; equivalently,
   \[
   f(g\cdot M)=f(M) \qquad \text{for all }g\in G_{\mathbf v} \text{ and }M\in\Rep_{\mathbf v}(\Pi).
   \]
   Set
   \[ \mathcal F(\Pi)
   =
   \bigoplus_{\mathbf v\in\mathbb N^I} \mathcal F_{\mathbf v}(\Pi).
   \]
   We regard these functions as functions on isomorphism classes of modules, extending each $f\in\mathcal F_{\mathbf v}(\Pi)$ by zero to modules whose dimension vectors differ from $\mathbf v$.

   Let $\chi_c$ denote the compactly supported Euler characteristic. For a constructible function $\varphi:X\to\mathbb C$, define its Euler integral by
   \[ \int_X \varphi\,d\chi_c
   =
   \sum_{c\in\Im(\varphi)} c\,\chi_c\bigl(\varphi^{-1}(c)\bigr).
   \]
   This sum is finite since $\varphi$ takes finitely many values.

   For a finite-dimensional $\Pi$-module $M$, let
   \[
   \Gr_{\mathbf v}^{\Pi}(M)
   =
   \{U\subseteq M: U \text{ is a }\Pi\text{-submodule and }
  \underline{\dim}U=\mathbf v\}.
   \]
   This is a closed subvariety of $\prod_{i\in I}\Gr_{v_i}(M_i)$, and hence is projective.

   For $f\in\mathcal F_{\mathbf v}(\Pi)$ and $g\in\mathcal F_{\mathbf w}(\Pi)$, define
   \begin{equation}\label{eq:Pi-convolution}
   (f*g)(M)
   =
   \int_{\Gr_{\mathbf v}^{\Pi}(M)} f(U)\,g(M/U)\,d\chi_c.
   \end{equation}
   $(\mathcal F(\Pi),*)$ is an $\mathbb N^I$-graded associative algebra. Its identity is the characteristic function $\mathbf 1_0$ of the zero module.
   
   For each $i\in I$, let $\theta_i$ be the characteristic function of the isomorphism class of the generalized simple module $E_i$:
   \[
   \theta_i(M)
   =
   \begin{cases}
   1,& M\cong E_i,\\
   0,& \text{otherwise}.
   \end{cases}
   \]
   Since $\underline{\dim}E_i=d_i\mathbf e_i$, we have $\theta_i\in\mathcal F_{d_i\mathbf e_i}(\Pi)$, where $(\mathbf e_i)_{i\in I}$ is the standard basis of $\mathbb Z^I$.

   \begin{definition}
   Let $\mathcal M(C,D)$ be the unital subalgebra of $(\mathcal F(\Pi),*)$ generated by the functions $\theta_i$:
   \[
   \mathcal M(C,D)
   =
   \mathbb C\langle\theta_i:i\in I\rangle_*.
   \]
   \end{definition}

   This is the algebra denoted by $\widetilde{\mathcal M}(\Pi)$ in \cite[Section~2.5]{GLS18}.

   We equip $\mathcal M(C,D)$ with the rank grading $\deg\theta_i=\mathbf e_i$. Equivalently, set
   \[
   \mathcal M_{\mathbf r}
   =
   \mathcal M(C,D)\cap\mathcal F_{D\mathbf r}(\Pi),
   \qquad
   \mathcal M(C,D)
   =
   \bigoplus_{\mathbf r\in\mathbb N^I}
   \mathcal M_{\mathbf r}.
   \]
   Each $\mathcal M_{\mathbf r}$ is finite-dimensional: it is spanned by the finitely many convolution words in which $\theta_i$ occurs $r_i$ times. Also, $\mathcal M_{\mathbf 0}=\mathbb C\mathbf 1_0$.

   For a sequence $\mathbf i=(i_1,\ldots,i_m)$, let $\Fl_{\mathbf i}^{E}(M)$ be the constructible set of flags of $\Pi$-submodules
   \[
   0=M^0\subset M^1\subset\cdots\subset M^m=M
   \]
   such that
   \[
   M^s/M^{s-1}\cong E_{i_s} \qquad (1\le s\le m).
   \]
   Iterating the convolution formula gives
   \begin{equation}\label{eq:Pi-flag-evaluation}
   (\theta_{i_1}*\cdots*\theta_{i_m})(M)
   =
   \chi_c\bigl(\Fl_{\mathbf i}^{E}(M)\bigr).
   \end{equation}

   Consequently, every element of $\mathcal M(C,D)$ vanishes on modules that are not $E$-filtered. The following is the analogue of \cite[Lemma~4.2]{GLS16} for $\Pi$. It is also recorded in \cite[Section~2.5]{GLS18}.
   \begin{lemma}\label{lem:locally-free-support}
   Let $X$ be a finite-dimensional $\Pi$-module. If $X$ is not locally free, then $f(X)=0$ for every $f\in\mathcal M(C,D)$.
   \end{lemma}
   
   We recall the following Hopf algebra structure from \cite[Section~2.5]{GLS18}; see also \cite[Propositions~4.5 and~4.7]{GLS16} for the corresponding construction for $H$.
   
   \begin{proposition}\label{prop:Hopf-structure}
   The convolution algebra $\mathcal M(C,D)$ is a connected cocommutative Hopf algebra. Its coproduct
   \[
   \Delta:\mathcal M(C,D)\longrightarrow \mathcal M(C,D)\otimes\mathcal M(C,D)
   \]
   is given by
   \[
   (\Delta f)(X,Y)=f(X\oplus Y).
   \]
   In particular,
   \[ \Delta(\theta_i) = \theta_i\otimes\mathbf1_0 + \mathbf1_0\otimes\theta_i.
   \]
   The counit is $\epsilon(f)=f(0)$, and the antipode satisfies $S(\theta_i)=-\theta_i$.
\end{proposition}
   Since $\Delta$ is an algebra homomorphism, for any sequence $\mathbf i=(i_1,\ldots,i_m)$ we have
   \[
   \Delta(\theta_{i_1}*\cdots*\theta_{i_m})
   =
   \prod_{a=1}^{m} \bigl(
   \theta_{i_a}\otimes\mathbf1_0 + \mathbf1_0\otimes\theta_{i_a} \bigr).
   \]

   \begin{definition}\label{def:primitive-element}
   Let $B$ be a Hopf algebra with coproduct $\Delta$ and unit $1_B$. An element $x\in B$ is called \emph{primitive} if
   \[
   \Delta(x)=x\otimes 1_B+1_B\otimes x.
   \]
   We denote the vector space of primitive elements by
   \[
   P(B)= \{x\in B:\Delta(x)=x\otimes 1_B+1_B\otimes x\}.
   \]
   \end{definition}
   
   The space $P(B)$ is a Lie subalgebra of $B$ under the commutator $[x,y]=xy-yx$. 
   %Indeed, since $\Delta$ is an algebra homomorphism, for $x,y\in P(B)$ we have
   %\[
   %\Delta([x,y])
   %=
   %[\Delta(x),\Delta(y)]
   %=
   %[x,y]\otimes 1_B+1_B\otimes[x,y].
   %\]
   Since each $\theta_i$ is primitive by Proposition~\ref{prop:Hopf-structure}, every iterated convolution commutator of the $\theta_i$ is also primitive.

   For the convolution algebra, primitive elements have the following interpretation; see also \cite[Section~2.5]{GLS18}.

   \begin{lemma}\label{lem:primitive-support}
   A function $f\in\mathcal M(C,D)$ is primitive if and only if
   \[
   f(0)=0 \qquad\text{and}\qquad f(X\oplus Y)=0 \quad\text{for all nonzero $\Pi$-modules $X,Y$}.
   \]
   Equivalently, a primitive function can take a nonzero value only on a nonzero indecomposable module.
\end{lemma}

\begin{proof}
    By Proposition~\ref{prop:Hopf-structure}, a function $f\in\mathcal M(C,D)$ is primitive if and only if
\[
   f(X\oplus Y)
   =
   f(X)\mathbf1_0(Y)+\mathbf1_0(X)f(Y)
\]
   for all finite-dimensional $\Pi$-modules $X,Y$. If $f$ is primitive, taking $X=Y=0$ gives $f(0)=0$, and taking $X,Y\ne0$ gives $f(X\oplus Y)=0$. Conversely, these two vanishing conditions imply the displayed identity: both sides vanish when $X,Y\ne0$, and the identity is immediate when one of the modules is zero.
\end{proof}
   As recalled in \cite[Section~2.5]{GLS18}, the inclusion
\[
  P(\mathcal M(C,D))\hookrightarrow\mathcal M(C,D)
\]
  induces a canonical Hopf algebra isomorphism
\[
  U\bigl(P(\mathcal M(C,D))\bigr) \xrightarrow{\sim} \mathcal M(C,D).
\]
Let $\mathfrak l(C,D)$ denote the Lie subalgebra of $\mathcal M(C,D)$ generated by the $\theta_i$ under the convolution commutator.
\begin{lemma}\label{lem:enveloping-realization}
    The inclusion $\mathfrak l(C,D)\hookrightarrow\mathcal M(C,D)$ induces a canonical Hopf algebra isomorphism
\[
    U(\mathfrak l(C,D))\xrightarrow{\sim}\mathcal M(C,D).
\]
\end{lemma}
\begin{proof}
    Since each $\theta_i$ is primitive, we have $\mathfrak l(C,D)\subseteq P(\mathcal M(C,D))$. By the PBW theorem, the induced map
    \[
    U(\mathfrak l(C,D))\longrightarrow\mathcal M(C,D)
\]
is injective. Its image contains all the algebra generators $\theta_i$, so it is surjective. 
\end{proof}

\subsection{The modified Cartan matrix and the main theorems}
\label{subsec:main-theorem}
Let $C=(c_{ij})_{i,j\in I}$ be a symmetrizable generalized Cartan matrix, and let $D=\operatorname{diag}(d_i)$ be a positive integral symmetrizer. Set
\[
    \lambda_i=
    \begin{cases}
        1, & d_i=1,\\
        2, & d_i>1,
    \end{cases}
\]
and define $A=A(C,D)=(a_{ij})_{i,j\in I}$ by
\begin{equation}\label{eq:modified-cartan}
    a_{ii}=2,\qquad
    a_{ij}=\lambda_i c_{ij}\quad(i\ne j).
\end{equation}
Then $A$ is a symmetrizable generalized Cartan matrix
with positive integral symmetrizer
\[
    D_A=\diag(2d_i/\lambda_i)_{i\in I}.
\]
We consider the condition
\begin{equation}\tag{E}\label{eq:edge-condition}
    \gcd(d_i,d_j)>1
    \qquad\text{whenever }c_{ij}<0
    \text{ and }d_i,d_j>1.
\end{equation}

Let $e_i$, $i\in I$, be the Chevalley generators of $\mathfrak n_+(A)$. We can now state the realization theorem.

\begin{theorem}\label{thm:main}
Assume that $(C,D)$ satisfies
condition~\eqref{eq:edge-condition}.
Then the assignment $e_i\mapsto\theta_i$ extends uniquely
to an isomorphism of rank-graded Hopf algebras
\[
    \Phi_{C,D}\colon
    U(\mathfrak n_+(A))
    \xrightarrow{\;\sim\;}
    \mathcal M(C,D).
\]
Its restriction gives a Lie algebra isomorphism
\[
    \mathfrak n_+(A)\xrightarrow{\;\sim\;} \mathfrak l(C,D).
\]
\end{theorem}

In Section~\ref{sec:surjectivity}, we prove the Serre relations associated with $A$ and establish surjectivity for arbitrary $(C,D)$, without assuming condition~\eqref{eq:edge-condition}. In Section~\ref{sec:injectivity}, we prove injectivity under this condition, completing the proof of Theorem~\ref{thm:main}.

\begin{corollary}[Finite and affine types] \label{cor:finite-affine}
Let $C$ be of finite or affine type, and let $D$ be any positive integral symmetrizer.  We have a canonical rank-graded Hopf algebra isomorphism
\[
\begin{aligned}
    U(\mathfrak n_+(A)) &\xrightarrow{\sim} \mathcal M(C,D),\\
    e_i &\longmapsto \theta_i,
\end{aligned}
\]
where $A=A(C,D)$. In particular, the Lie subalgebra generated by the $\theta_i$ under the convolution commutator is canonically isomorphic to $\mathfrak n_+(A)$.
\end{corollary}

\begin{proof}
Since $C$ is of finite or affine type, the symmetric matrix $DC$ is positive semidefinite. Its two-by-two principal minors therefore
give
\[
    c_{ij}c_{ji}\leq 4
    \qquad (i\neq j).
\]
It is then easy to check that condition~\eqref{eq:edge-condition} holds. The assertions now follow from Theorem~\ref{thm:main}.
\end{proof}

For any symmetrizable generalized Cartan matrix $C$ and any positive integral symmetrizer $D'$, the pair $(C,mD')$ satisfies condition~\eqref{eq:edge-condition} for every integer $m\ge 2$, since all diagonal entries of $mD'$ are divisible by $m$. Thus Theorem~\ref{thm:main} applies to $(C,mD')$.

For pairs $(C,D)$ not satisfying condition~\eqref{eq:edge-condition}, we prove the following theorem in Section~\ref{app:serre-quotients}.

\begin{theorem}\label{thm:nonrealization}
   If $(C,D)$ does not satisfy condition~\eqref{eq:edge-condition}, then
\[
    \mathcal M(C,D)\not\cong U(\mathfrak n_+(B))
\]
as unital $\mathbb C$-algebras for any symmetrizable generalized Cartan matrix $B$. In particular, $\mathfrak l(C,D)$ is not isomorphic to the positive part of any symmetrizable Kac--Moody Lie algebra. 
\end{theorem}



%%%%%%%%%%%%%%%%%%%%%%%%%%%%%%%%%%%%%%%%%%%%%55
\section{Serre relations and the canonical homomorphism} \label{sec:surjectivity}
In this section, we prove the Serre relations associated with $A=A(C,D)$ and obtain the canonical homomorphism to $\M(C,D)$. We do not assume condition~\eqref{eq:edge-condition}.
%let $C$ be a symmetrizable generalized Cartan matrix and $D$ a positive integral symmetrizer, and let $A=A(C,D)$ be defined by \eqref{eq:modified-cartan}. We prove that the generators $\theta_i$ satisfy the Serre relations associated with $A$, yielding a canonical surjective Hopf algebra homomorphism \[ U(\mathfrak n_+(A))\twoheadrightarrow\mathcal M(C,D). \] No assumption of condition~\eqref{eq:edge-condition} is needed.
We use the incoming and outgoing maps from \cite[Sections~5.1--5.3]{GLS17}. 

For $c_{ij}<0$, let ${}_iH_j$ denote the $(H_i,H_j)$-bimodule of arrows from $j$ to $i$. It is free of rank $|c_{ij}|$ as a left $H_i$-module. For a $\Pi$-module $M$, set
\[
    \widetilde M_i
    =
    \bigoplus_{j:\,c_{ij}<0}
    {}_iH_j\otimes_{H_j}M_j.
\]
Following \cite[Section~5.3]{GLS17}, we consider the $H_i$-module homomorphisms
\[
    M_{i,\mathrm{in}}\colon\widetilde M_i\longrightarrow M_i,
    \qquad
    M_{i,\mathrm{out}}\colon M_i\longrightarrow\widetilde M_i.
\]
By \cite[Proposition~5.2]{GLS17}, the preprojective relations are equivalent to
\[
    M_{i,\mathrm{in}}\cdot M_{i,\mathrm{out}}=0
    \qquad(i\in I).
\]



\subsection{Splitting lemma}\label{subsec:splitting}
We adapt the splitting argument in the proof of \cite[Lemma~6.3]{GLS18} to locally free $\Pi$-modules.
\begin{lemma}\label{lem:general splitting}
   Let $X$ be a locally free $\Pi$-module and let $i\in I$. Let $m$ and $n$ be the minimal numbers of $H_i$-generators of $\Im X_{i,\mathrm{in}}$ and $\Im X_{i,\mathrm{out}}$. If $\rank_{H_i}X_i>m+n$, then $E_i$ is a direct summand of $X$.
\end{lemma}

\begin{proof}
    Write $\alpha=X_{i,\mathrm{in}}$ and $\beta=X_{i,\mathrm{out}}$. Thus
\[
    \widetilde X_i \xrightarrow{\alpha} X_i
    \xrightarrow{\beta} \widetilde X_i,
    \qquad \alpha\beta=0.
\]
We put the matrix of $\alpha$ in Smith normal form. Since $\Im \alpha$ is minimally generated by $m$ elements, this form has exactly $m$ nonzero diagonal entries. Thus we obtain a decomposition into free $H_i$-modules
\[
    X_i=P\oplus Z,\qquad
    P\cong H_i^m,\qquad
    \Im \alpha\subseteq P.
\]
Similarly, applying Smith normal form to $\beta$, we obtain a decomposition into free $H_i$-modules
\[
    \widetilde X_i=T\oplus T',\qquad
    T\cong H_i^n,\qquad
    \Im \beta\subseteq T.
\]
Now consider the restriction $\beta|_Z:Z\to T$. Since
\[
    \rank_{H_i}Z
    =\rank_{H_i}X_i-m>n=\rank_{H_i}T,
\]
we may choose a direct summand $K\cong H_i$ of $Z$ such that $\beta(K)=0$. Since $\Im \alpha\subseteq P$, $E_i$ is a direct summand of $X$.
\end{proof}

\begin{corollary}\label{lem:splitting}
Let $i\ne j$ with $c_{ij}<0$, and let $X$ be a locally free $\Pi$-module of rank $r\mathbf e_i+\mathbf e_j$. If $r>\lambda_i|c_{ij}|$, then $E_i$ is a direct summand of $X$.
\end{corollary}

\begin{proof}
    Put $N=|c_{ij}|$. Since $X$ is locally free, we have
\begin{align*}
    X_i\cong H_{i}^r,\qquad X_j\cong H_j,\qquad
    \widetilde X_i
    ={}_iH_j\otimes_{H_j}X_j\cong H_i^N.
\end{align*}
Write $\alpha=X_{i,\mathrm{in}}$ and $\beta=X_{i,\mathrm{out}}$. Let $m,n$ be as in Lemma~\ref{lem:general splitting}.

If $d_i=1$, then $H_i=\mathbb C$. Since
$\operatorname{Im}\beta\subseteq\ker\alpha$, we have
\[
    m+n
    =\operatorname{rank}\alpha+\operatorname{rank}\beta
    \leq N.
\]
If $d_i>1$, we have
\[
    H_i^N \xrightarrow{\alpha} H_i^r \xrightarrow{\beta} H_i^N.
\]
Thus $m \leq N$ and $n \leq N$.

In both cases, $m+n\leq\lambda_i N<r$. Thus $E_i$ is a direct summand of $X$ by Lemma~\ref{lem:general splitting}.
\end{proof}

\subsection{The Serre relations}\label{subsec:Serre}
We now apply Corollary~\ref{lem:splitting} to prove the Serre relations associated with $A(C,D)$. For $f,g\in\mathcal M(C,D)$, write
\[
   \ad(f)(g)=[f,g]=f*g-g*f.
\]

\begin{proposition}\label{prop:A-Serre}
The generators $\theta_i$, $i\in I$, satisfy
\begin{equation}\label{eq:A-Serre}
(\ad\theta_i)^{1-a_{ij}}(\theta_j)=0 \qquad (i\ne j).
\end{equation}
\end{proposition}

\begin{proof}
    Fix $i \neq j$ and put
\[
   r=1-a_{ij}=1+\lambda_i|c_{ij}|, \qquad \Theta_{ij}=(\ad\theta_i)^r(\theta_j).
\]
$\Theta_{ij}$ is primitive and homogeneous of rank degree $r\mathbf e_i+\mathbf e_j$. By Lemma~\ref{lem:locally-free-support} and Lemma~\ref{lem:primitive-support}, it can be nonzero only on locally free indecomposable modules of this rank.

Let $X$ be a locally free module with rank $r\mathbf e_i+\mathbf e_j$. If $c_{ij}=0$, then $r=1$ and $X \cong E_i \oplus E_j$. If $c_{ij}<0$, Corollary ~\ref{lem:splitting} gives $X\cong E_i\oplus Y$. Thus each such $X$ is decomposable and $\Theta_{ij}=0$.
\end{proof}

\begin{corollary}\label{cor:canonical-surjection}
For every symmetrizable generalized Cartan matrix $C$ and every positive integral symmetrizer $D$, the assignment $e_i\mapsto\theta_i$ extends uniquely to a surjective rank-graded Hopf algebra homomorphism
\begin{equation}\label{eq:canonical-surjection}
\Phi_{C,D}\colon U\bigl(\nplus(A(C,D))\bigr) \twoheadrightarrow \mathcal M(C,D).
\end{equation}
\end{corollary}

\begin{proof}
 The Serre presentation and Proposition~\ref{prop:A-Serre} give the unique algebra homomorphism with $e_i\mapsto\theta_i$. It is surjective because $\theta_i$ generate $\M(C,D)$ and it preserves the grading and the Hopf structure because the generators on both sides are primitive of degree $\mathbf e_i$.
\end{proof}






%%%%%%%%%%%%%%%%%%%%%%%%%%%%%%%%%%
\section{Injectivity of the canonical homomorphism} \label{sec:injectivity}

Put $A=A(C,D)$. We first construct an auxiliary quiver $Q$ for an arbitrary pair $(C,D)$. Under condition~\eqref{eq:edge-condition}, 
each representation $V$ of $Q$ gives a locally free $\Pi(C,D)$-module $F(V)$. To prove that $\Phi_{C,D}$ is injective, we use a torus action to compare the $E$-flags of $F(V)$ with flags of $Q$-subrepresentations of $V$. This comparison is closely related to the argument in \cite[Appendix~A.2]{Fei26}. In the construction of $F(V)$, we choose weight shifts for the arrow maps and use condition~\eqref{eq:edge-condition} to verify the mesh relations.

\subsection{Quiver convolution algebras and symmetric coverings}
Let $Q$ be a finite quiver without loops, with vertex set $Q_0$; oriented cycles and multiple arrows are allowed. Let $B=(b_{ij})_{i,j\in Q_0}$ be its associated symmetric generalized Cartan matrix:
\[
  b_{ii}=2,\qquad b_{ij} = -\#\{\text{arrows joining $i$ and $j$ in either direction}\} \quad(i\ne j).
\]

For $\mathbf d\in\mathbb N^{Q_0}$, let $\mathcal F_{\mathbf d}(Q)$ be the space of
isomorphism-invariant $\mathbb C$-valued constructible functions on $\Rep_{\mathbf d}(Q)$. The space
\[
  \mathcal F(Q) = \bigoplus_{\mathbf d\in\mathbb N^{Q_0}} \mathcal F_{\mathbf d}(Q)
\]
  is equipped with the convolution product
\[
  (f*g)(V) = \int_{W\subseteq V} f(W)g(V/W)\,d\chi_c,
\]
where $W$ runs through the subrepresentations of $V$.

For each $i\in Q_0$, let $S_i$ be the one-dimensional representation supported at $i$, and put $s_i=\mathbf1_{S_i}$. We denote by
\[
  \mathcal M(Q) = \mathbb C\langle s_i:i\in Q_0\rangle_* \subseteq\mathcal F(Q)
\]
the unital convolution subalgebra generated by the $s_i$.

We recall the classical constructible-function realization of the positive enveloping algebra, originating in the work of Schofield and Lusztig \cite{Sch91,Lus91}. For an explicit statement for quivers without loops, see \cite[Theorem~3.1]{JS12}.

\begin{theorem}\label{thm:inj-hall}
The assignment $e_i\mapsto s_i$ extends uniquely
to a graded algebra isomorphism
\[
\eta_Q\colon U(\nplus(B))
\xrightarrow{\sim}\mathcal M(Q).
\]
\end{theorem}

%In \cite{JS12}, the convolution algebra is defined on nilpotent representations, with the opposite convolution convention. Every convolution word in the $s_i$ is supported on representations admitting a filtration with vertex-simple factors, and hence on the nilpotent locus. Restriction to this locus therefore identifies $\mathcal M(Q)$ with the opposite of the composition algebra in \cite{JS12}. The stated isomorphism then follows by composing with the algebra anti-automorphism of $U(\mathfrak n_+(B))$ that fixes each $e_i$ and reverses the order of products. For a formulation using constructible functions on the full representation spaces, see also \cite[\S9.2]{Hen24}.

   For a sequence $\mathbf j=(j_1,\ldots,j_r)$ of vertices of $Q$ and a finite-dimensional representation $V$, let $\Fl^{Q}_{\mathbf j}(V)$ denote the variety of flags of subrepresentations
\[
   0=V^0\subset V^1\subset\cdots\subset V^r=V
\]
  such that
\[
  V^k/V^{k-1}\simeq S_{j_k} \qquad (1\leq k\leq r).
\]
  With our convolution convention, we have
\begin{equation}\label{eq:inj-Q-flags}
  (s_{j_1}*\cdots*s_{j_r})(V) = \chi_c\bigl(\Fl^{Q}_{\mathbf j}(V)\bigr).
\end{equation}
For integers $m>0$ and $w$, write $[w]_{m}$ for the least nonnegative residue of $w$ modulo $m$. The following folding lemma is due to Hubery \cite[Theorem~7.1.5]{Hub02}; see also \cite[Proposition~4.4]{Dem11}.

\begin{lemma}\label{lem:inj-cover}
Let $A=(a_{ij})_{i,j\in I}$ be a symmetrizable generalized Cartan matrix, and let $B$ be a symmetric generalized Cartan matrix indexed by
\[
\{(i,a):i\in I,\ 0\leq a<m_i\}.
\]
Here each $m_i$ is a positive integer. Assume that the permutation $\sigma(i,a)=(i,[a+1]_{m_i})$ preserves $B$, that no two distinct vertices over the same $i$ are adjacent, and that
\begin{equation}\label{eq:inj-columnsum}
\sum_{a=0}^{m_i-1}b_{(i,a),(j,c)}
=
a_{ij}
\qquad
(i,j\in I,\ 0\leq c<m_j).
\end{equation}
Then $e_i\longmapsto \sum_{a=0}^{m_i-1}e_{(i,a)}$ defines an injective Lie algebra homomorphism
\[
\iota\colon
\nplus(A)\hookrightarrow\nplus(B).
\]
The induced homomorphism $U(\iota)$ is also injective.
\end{lemma}

\begin{proof}
The cyclic group $\langle\sigma\rangle$ acts admissibly on the diagram of $B$, and \eqref{eq:inj-columnsum} identifies its folding Cartan matrix with $A$. The Lie algebra embedding is therefore given by \cite[Proposition~4.4]{Dem11}, and PBW gives the assertion about enveloping algebras.
\end{proof}

\begin{remark}
When $B$ comes from a quiver, $\sigma$ need only preserve the underlying graph, not its arrow directions.
\end{remark}




%%%%%%%%%%%%%%%%%%%%%%%%%%%%%%%%%%%%%%%%%%%%%%%%%
\subsection{Construction of the auxiliary quiver and the folding embedding}\label{subsec:auxiliary-quiver}
  We construct a finite quiver $Q$ without loops whose associated symmetric generalized Cartan matrix $B$ folds to $A=A(C,D)$. The construction and the embedding below do not require condition~\eqref{eq:edge-condition}.

\subsubsection{Vertices}
Set
\begin{equation}\label{eq:inj-weightsizes}
  L=\lcm\bigl(2,d_i:i\in I\bigr), \qquad
  m_i=
\begin{cases}
  L/2,& d_i=1,\\
  L/d_i,& d_i>1.
\end{cases}
\end{equation}
  Define
\[
    Q_0=\{(i,a):i\in I,\ 0\le a<m_i\},
    \qquad
    I_i=\{(i,a):0\le a<m_i\}.
\]

\subsubsection{Arrows}
For each unordered adjacent pair $\{i,j\}$ of $C$, apply exactly one of the following rules. In each case, $1\le\ell\le g_{ij}$.

\par\smallskip
\noindent\textbf{Case 1: $d_i=d_j=1$.}
Here $m_i=m_j=L/2$. Choose and fix a direction, say $j\to i$. For each $0\leq w<L/2$
introduce an arrow
\[
\alpha_{ij,w}^{(\ell)}:
(j,w)\longrightarrow(i,w).
\]

\par\smallskip
\noindent\textbf{Case 2: exactly one of $d_i,d_j$ equals $1$.}
After interchanging $i$ and $j$ if necessary, assume that $d_i>1$ and $d_j=1$. Here $m_i=L/d_i$ and $m_j=L/2$. For each $0\le w<L/2$, introduce the arrows
\begin{equation}\label{eq:inj-mixed-arrows}
\begin{aligned}
\beta_{ij,w}^{(\ell)}:
(i,[w]_{m_i})&\longrightarrow(j,w),\\
\alpha_{ij,w}^{(\ell)}:
(j,w)&\longrightarrow(i,[w+L/2]_{m_i}).
\end{aligned}
\end{equation}

\par\smallskip
\noindent\textbf{Case 3: $d_i,d_j>1$.}
Here $m_i=L/d_i$ and $m_j=L/d_j$. Recall that $k_{ij}=\gcd(d_i,d_j)$, and put
\[
t=L/k_{ij}=\lcm(m_i,m_j).
\]
For each $0\le w<t$, introduce the arrows
\begin{equation}\label{eq:auxiliary-long-arrows}
\begin{aligned}
    \alpha^{(\ell)}_{ij,w}&\colon
    (j,[w]_{m_j})\longrightarrow(i,[w]_{m_i}),\\
    \beta^{(\ell)}_{ij,w}&\colon
    (i,[w]_{m_i})\longrightarrow(j,[w]_{m_j}).
\end{aligned}
\end{equation}
There are no other arrows.

\subsubsection{The folding embedding}

Let
\[
    B=\bigl(b_{(i,a),(j,b)}\bigr)
\]
be the symmetric generalized Cartan matrix associated with $Q$.
For each $i\in I$, set
\[
    \sigma_i=\sum_{a=0}^{m_i-1}s_{(i,a)}
    \in\mathcal M(Q).
\]

\begin{theorem}\label{thm:auxiliary-embedding}
The assignment
\[
    e_i\longmapsto\sum_{a=0}^{m_i-1}e_{(i,a)}
\]
extends to an injective Lie algebra homomorphism
\[
    \iota\colon
    \mathfrak n_+(A)\hookrightarrow\mathfrak n_+(B).
\]
Consequently, the assignment $e_i\mapsto\sigma_i$
extends uniquely to an injective algebra homomorphism
\[
    J=\eta_Q\circ U(\iota)\colon
    U(\mathfrak n_+(A))\hookrightarrow\mathcal M(Q).
\]
\end{theorem}

\begin{proof}
    No arrow joins two vertices in the same fiber $I_i$. We check that $\sigma(i,a)=(i,[a+1]_{m_i})$ preserves $B$ and that
    \begin{align}
        \sum_{a=0}^{m_i-1}b_{(i,a),(j,b)}=a_{ij}.
    \end{align}
The identity is immediate when $i=j$ or $c_{ij}=0$. We therefore assume that $c_{ij}<0$ and consider the following three cases.
\par\smallskip
\noindent\textbf{Case 1: $d_i=d_j=1$.}
Here $m_i=m_j=L/2$. 
For each $0\le w<L/2$, there are $g_{ij}$ edges between $(i,w)$ and $(j,w)$. $\sigma$ preserves these edges and
\[
    \sum_{a=0}^{m_i-1}b_{(i,a),(j,b)}
    =-g_{ij}=c_{ij}=a_{ij}.
\]
The same argument applies with $i$ and $j$ interchanged.

\par\smallskip
\noindent\textbf{Case 2: exactly one of $d_i,d_j$ equals $1$.} After interchanging $i$ and $j$ if necessary, assume that $d_i>1$ and $d_j=1$. Then $m_i=L/d_i$ and $m_j=L/2$, and
\[
    c_{ij}=-g_{ij},
    \qquad
    c_{ji}=-d_i g_{ij}.
\]

For each $0\le b<m_j$, the vertex $(j,b)$ is joined to $(i,[b]_{m_i})$ and $(i,[b+m_j]_{m_i})$ by $g_{ij}$ edges each, counted with multiplicity. After forgetting the arrow directions, $\sigma$ sends each edge in either family at $b$ to the corresponding edge in the same family at $b+1$, provided $b<m_j-1$. At $b=m_j-1$, it sends the edge underlying $\beta^{(\ell)}_{ij,m_j-1}$ to that underlying $\alpha^{(\ell)}_{ij,0}$, and the edge underlying $\alpha^{(\ell)}_{ij,m_j-1}$ to that underlying $\beta^{(\ell)}_{ij,0}$, since $2m_j=L$ is divisible by $m_i$. Thus $\sigma$ preserves the edge multiplicities between the two fibers.

Each vertex $(j,b)$ is incident with $2g_{ij}$ edges joining it to $I_i$, so
\[
    \sum_{a=0}^{m_i-1}b_{(i,a),(j,b)}
    =-2g_{ij}
    =2c_{ij}
    =a_{ij}.
\]
The indices $b$ and $b+m_j$, as $b$ ranges from $0$ to $m_j-1$, together run through $0,\ldots,L-1$. Each residue modulo $m_i$ therefore occurs $L/m_i=d_i$ times. Hence
\[
    \sum_{b=0}^{m_j-1}b_{(j,b),(i,a)}
    =-d_i g_{ij}
    =c_{ji}
    =a_{ji}.
\]

\par\smallskip
\noindent\textbf{Case 3: $d_i,d_j>1$.}
Put
\[
    t=\frac{L}{k_{ij}}
      =\lcm(m_i,m_j).
\]
Each index $0\le w<t$ contributes $2g_{ij}$ edges between
$(i,[w]_{m_i})$ and $(j,[w]_{m_j})$. Since $m_i$ and $m_j$ divide $t$, $\sigma$ corresponds to replacing $w$ by $[w+1]_t$ and preserves the edge multiplicities.

For each fixed $0\le b<m_j$, exactly $t/m_j$ values of $w$ satisfy $[w]_{m_j}=b$. Since
\[
    \frac{t}{m_j}
    =\frac{d_j}{k_{ij}}
    =f_{ij},
\]
we obtain
\[
    \sum_{a=0}^{m_i-1}b_{(i,a),(j,b)}
    =-2g_{ij}\frac{t}{m_j}
    =-2g_{ij}f_{ij}
    =2c_{ij}
    =a_{ij}.
\]
Similarly, $t/m_i=d_i/k_{ij}=f_{ji}$, and hence
\[
    \sum_{b=0}^{m_j-1}b_{(j,b),(i,a)}
    =-2g_{ij}\frac{t}{m_i}
    =-2g_{ij}f_{ji}
    =2c_{ji}
    =a_{ji}.
\]
The result now follows from Lemma~\ref{lem:inj-cover} and Theorem~\ref{thm:inj-hall}.
\end{proof}

\begin{remark}
After forgetting the arrow directions, the construction of $Q$ admits a uniform description. For each unordered adjacent pair $\{i,j\}$, set
\[
    t_{ij}=\lcm(m_i,m_j),
    \qquad
    \rho_{ij}=\gcd(|a_{ij}|,|a_{ji}|).
\]
For each $0\le w<t_{ij}$, join $(i,[w]_{m_i})$ and $(j,[w]_{m_j})$ by $\rho_{ij}$ edges. Then $\sigma$ preserves the edge multiplicities, and
\[
    \rho_{ij}\frac{t_{ij}}{m_j}=|a_{ij}|,
    \qquad
    \rho_{ij}\frac{t_{ij}}{m_i}=|a_{ji}|.
\]
This gives the folding relations. We use the chosen arrow directions to construct locally free $\Pi(C,D)$-modules from $Q$-representations in the next subsection.
\end{remark}

\begin{example}\label{ex:auxiliary-quiver}
Let
\[
    D=\diag(1,1,2,4),
    \qquad
    C=
    \begin{pmatrix}
        2&-1&0&0\\
        -1&2&-2&0\\
        0&-1&2&-2\\
        0&0&-1&2
    \end{pmatrix}.
\]
Then $DC$ is symmetric, and
\[
    A(C,D)=
    \begin{pmatrix}
        2&-1&0&0\\
        -1&2&-2&0\\
        0&-2&2&-4\\
        0&0&-2&2
    \end{pmatrix}.
\]
We have $g_{12}=g_{23}=g_{34}=1$ and
\[
    L=4,
    \qquad
    (m_1,m_2,m_3,m_4)=(2,2,2,1).
\]
Thus
\[
    Q_0=
    \{(i,a):1\le i\le3,\ a=0,1\}
    \sqcup\{(4,0)\}.
\]
Choosing the direction $2\to1$ in Case~1, the resulting quiver is
\[
    (1,0)
    \xleftarrow{\alpha_{12,0}}
    (2,0)
    \overset{\alpha_{32,0}}
      {\underset{\beta_{32,0}}{\rightleftarrows}}
    (3,0)
    \overset{\beta_{34,0}}
      {\underset{\alpha_{34,0}}{\rightleftarrows}}
    (4,0)
    \overset{\alpha_{34,1}}
      {\underset{\beta_{34,1}}{\rightleftarrows}}
    (3,1)
    \overset{\beta_{32,1}}
      {\underset{\alpha_{32,1}}{\rightleftarrows}}
    (2,1)
    \xrightarrow{\alpha_{12,1}}
    (1,1).
\]
\end{example}

\subsection{From quiver representations to
generalized preprojective modules}
Assume condition~\eqref{eq:edge-condition}. Let $V$ be a finite-dimensional representation of the auxiliary quiver $Q$. Write $V_{i,a}$ for its space at $(i,a)$ and $V(\rho)$ for its linear map along an arrow $\rho$. We construct a locally free $\Pi(C,D)$-module $F(V)$
by specifying its vertex spaces and arrow actions.

\subsubsection{Vertex spaces and weights}
For each $i\in I$, define
\begin{equation}\label{eq:inj-testvertices}
F(V)_i
=
\bigoplus_{a=0}^{m_i-1}
H_i\otimes_{\mathbb C}V_{i,a}.
\end{equation}
The loop $\varepsilon_i$ acts by multiplication
on the $H_i$-factor.

For $0\le a<m_i$ and $0\le p<d_i$, assign weight $a+m_ip$ to the subspace
\[
    F_i[a+m_ip]
    :=
    \mathbb C\varepsilon_i^p
    \otimes_{\mathbb C}V_{i,a}.
\]
The admissible weights are
\[
    \begin{cases}
        0\le w<L/2, & d_i=1,\\
        0\le w<L,   & d_i>1.
    \end{cases}
\]
For all other integers $w$, set $F_i[w]=0$. For every admissible weight $w=a+m_ip$, we use
the $\mathbb C$-linear isomorphism
\[
    F_i[w]\xrightarrow{\;\sim\;}V_{i,[w]_{m_i}},
    \qquad
    \varepsilon_i^p\otimes v\longmapsto v.
\]
The loop action satisfies
\[
\varepsilon_i F_i[w]
\subseteq F_i[w+m_i].
\]

\subsubsection{Arrow actions}

For each unordered adjacent pair $\{i,j\}$,
use the same labeling as in the construction
of $Q$.
For $1\leq\ell\leq g_{ij}$, we define maps
\[
a_\ell:F(V)_j\longrightarrow F(V)_i,
\qquad
b_\ell:F(V)_i\longrightarrow F(V)_j
\]
as the actions of the arrows
$\alpha_{ij}^{(\ell)}$ and
$\alpha_{ji}^{(\ell)}$ of $\Pi(C,D)$,
respectively.

The maps on weight spaces below are defined via the canonical $\mathbb C$-linear identifications
\[
    F_i[w]\cong V_{i,[w]_{m_i}}.
\]
All unlisted weight blocks are zero.

\par\smallskip
\noindent\textbf{Case 1: $d_i=d_j=1$.}
Retain the direction $j\to i$ chosen in the construction of $Q$.
For $0\leq w<L/2$, define
\[
a_\ell\big|_{F_j[w]}:
F_j[w]
\xrightarrow{\;V(\alpha_{ij,w}^{(\ell)})\;}
F_i[w],
\]
and set $b_\ell=0$.

\par\smallskip
\noindent\textbf{Case 2: $d_i>1$ and $d_j=1$.}
For $0\leq w<L/2$, define
\begin{equation}\label{eq:inj-mixedplacement}
\begin{aligned}
b_\ell\big|_{F_i[w]}:\quad
F_i[w]
&\xrightarrow{\;V(\beta_{ij,w}^{(\ell)})\;}
F_j[w],\\
a_\ell\big|_{F_j[w]}:\quad
F_j[w]
&\xrightarrow{\;V(\alpha_{ij,w}^{(\ell)})\;}
F_i[w+L/2].
\end{aligned}
\end{equation}
In particular, $b_\ell$ is zero on all weight spaces of weight at least $L/2$.

\par\smallskip
\noindent\textbf{Case 3: $d_i,d_j>1$.}
Put
\[
t=\frac{L}{k_{ij}},
\qquad
s=L-t,
\]
where $k_{ij}=\gcd(d_i,d_j)$. For $0\leq w<t$, define
\begin{equation}\label{eq:inj-thickplacement}
\begin{aligned}
a_\ell\big|_{F_j[w]}:\quad
F_j[w]
&\xrightarrow{\;V(\alpha_{ij,w}^{(\ell)})\;}
F_i[w+s],\\
b_\ell\big|_{F_i[w]}:\quad
F_i[w]
&\xrightarrow{\;V(\beta_{ij,w}^{(\ell)})\;}
F_j[w+s].
\end{aligned}
\end{equation}
Both maps are zero on source weight spaces of weight at least $t$.

Since $t=\lcm(m_i,m_j)$ and both
$m_i$ and $m_j$ divide $L$, the shift $s=L-t$
is divisible by both $m_i$ and $m_j$.
Consequently,
\[
    [w+s]_{m_i}=[w]_{m_i},
    \qquad
    [w+s]_{m_j}=[w]_{m_j}.
\]
Thus the source and target weight spaces identify with the source and target spaces of the corresponding arrow maps of $V$.

\begin{lemma}\label{lem:relations}
Assume condition~\eqref{eq:edge-condition}. For every finite-dimensional representation $V$ of $Q$, the operators defined above endow
\[
    F(V)=\bigoplus_{i\in I}F(V)_i
\]
with the structure of a locally free $\Pi(C,D)$-module.
Moreover,
\[
    \rank_{H_i}F(V)_i
    =
    \sum_{a=0}^{m_i-1}\dim_{\mathbb C}V_{i,a}
    \qquad(i\in I).
\]
\end{lemma}

\begin{proof}
The relations $\varepsilon_i^{d_i}=0$ hold by the definition of $H_i$.
We prove the stronger assertion that both sides of every commutativity relation and each individual summand in the mesh relations act as zero on $F(V)$. We check this for each adjacent pair $\{i,j\}$ and each $1\le\ell\le g_{ij}$.
\par\smallskip
\noindent\textbf{Case 1: $d_i=d_j=1$.}
The arrows of $Q$ between $I_i$ and $I_j$ all have the chosen direction $j\to i$. For the reverse arrow $\alpha_{ji}^{(\ell)}$ of $\Pi(C,D)$, we set $b_\ell=0$.
Since $d_i=d_j=1$, both loops act as zero, so both sides of the commutativity relations vanish. Moreover, $b_\ell=0$ implies
\[
    a_\ell b_\ell=b_\ell a_\ell=0.
\]

\par\smallskip
\noindent\textbf{Case 2: $d_i>1$ and $d_j=1$.}
Here $f_{ij}=1$ and $f_{ji}=d_i$.
Since $\varepsilon_i^{d_i}=0$ and $\varepsilon_j=0$, both sides of the commutativity relations vanish:
\[
    \varepsilon_i^{d_i}a_\ell
    =0=a_\ell\varepsilon_j,
    \qquad
    \varepsilon_jb_\ell
    =0=b_\ell\varepsilon_i^{d_i}.
\]

The image of $a_\ell$ lies in weights
at least $L/2$, on which $b_\ell$ is zero.
Thus
\[
b_\ell a_\ell=0,
\]
which gives the contribution to the mesh relation at $j$.

At $i$, the corresponding mesh terms are
\[
\varepsilon_i^p a_\ell b_\ell
\varepsilon_i^{d_i-1-p},
\qquad 0\leq p<d_i.
\]
The map $a_\ell b_\ell$ raises weight by
$L/2$, while each loop raises weight by
$m_i$. Hence each displayed term raises
weight by
\[
\frac L2+(d_i-1)m_i
=
\frac{3L}{2}-m_i
\geq L,
\]
where we used $d_i\geq2$ and $m_i=L/d_i$. Since $F_i[w]=0$ outside $0\le w<L$, every such mesh term is zero.

\par\smallskip
\noindent\textbf{Case 3: $d_i,d_j>1$.}
Recall that
\[
    t=\frac{L}{k_{ij}},
    \qquad
    s=L-t,
\]
and
\[
    d_i=k_{ij}f_{ji},
    \qquad
    d_j=k_{ij}f_{ij}.
\]
Hence
\[
    f_{ji}m_i=f_{ij}m_j=t,
\]
so the loop powers $\varepsilon_i^{f_{ji}}$ and $\varepsilon_j^{f_{ij}}$ both raise weight by $t$.

Both $a_\ell$ and $b_\ell$ raise weight by $s$, while $\varepsilon_i^{f_{ji}}$ and $\varepsilon_j^{f_{ij}}$ raise weight by $t$. Hence each side of the commutativity relations raises weight by $s+t=L$. Since all admissible weights at $i$ and $j$ lie between $0$ and $L-1$, we obtain
\[
    \varepsilon_i^{f_{ji}}a_\ell
    =a_\ell\varepsilon_j^{f_{ij}}=0,
    \qquad
    \varepsilon_j^{f_{ij}}b_\ell
    =b_\ell\varepsilon_i^{f_{ji}}=0.
\]

At vertex $i$, consider a mesh term
\[
    \varepsilon_i^p a_\ell b_\ell
    \varepsilon_i^{f_{ji}-1-p},
    \qquad 0\le p<f_{ji}.
\]
Each arrow map raises weight by $s$, while each application of $\varepsilon_i$ raises
weight by $m_i$. Hence this term raises weight by
\[
    2s+(f_{ji}-1)m_i
    =
    2(L-t)+t-m_i
    =
    2L-t-m_i.
\]
Under condition~\eqref{eq:edge-condition}, we have $k_{ij}\ge2$, so $t=L/k_{ij}\le L/2$. Moreover, $m_i\le t$ since $t=f_{ji}m_i$. Consequently,
\[
    2L-t-m_i\ge 2L-2t\ge L.
\]
Since $F_i[w]=0$ outside $0\le w<L$, the displayed mesh term is zero. Interchanging $i$ and $j$ gives the same conclusion for every mesh term at $j$.

Summing these zero terms with the prescribed signs proves the mesh relations at every vertex. Together with the loop and commutativity relations, this shows that the operators define a $\Pi(C,D)$-module structure on $F(V)$.

Finally, each summand $H_i\otimes_{\mathbb C}V_{i,a}$ is free over $H_i$ of rank $\dim_{\mathbb C}V_{i,a}$. Thus $F(V)$ is locally free, with
\[
    \rank_{H_i}F(V)_i
    =
    \sum_{a=0}^{m_i-1}\dim_{\mathbb C}V_{i,a}
    \qquad(i\in I).
\]
\end{proof}

\begin{example}\label{ex:auxiliary-module}
Continue with Example~\ref{ex:auxiliary-quiver}.
Take $V_{i,a}=\mathbb C$ at every vertex of $Q$ and assign
arbitrary scalars $x_a,y_a,z_a,u_a,v_a\in\mathbb C$,
for $a=0,1$, to its arrows as follows:
\[
\begin{tikzcd}[column sep=1.5em]
    \underset{(1,0)}{\mathbb C}
    & \underset{(2,0)}{\mathbb C}
        \arrow[l,"x_0"']
        \arrow[r,shift left=.6ex,"y_0"]
    & \underset{(3,0)}{\mathbb C}
        \arrow[l,shift left=.6ex,"z_0"]
        \arrow[r,shift left=.6ex,"v_0"]
    & \underset{(4,0)}{\mathbb C}
        \arrow[l,shift left=.6ex,"u_0"]
        \arrow[r,shift left=.6ex,"u_1"]
    & \underset{(3,1)}{\mathbb C}
        \arrow[l,shift left=.6ex,"v_1"]
        \arrow[r,shift left=.6ex,"z_1"]
    & \underset{(2,1)}{\mathbb C}
        \arrow[l,shift left=.6ex,"y_1"]
        \arrow[r,"x_1"]
    & \underset{(1,1)}{\mathbb C}
\end{tikzcd}
\]
Here each arrow label denotes multiplication by the
indicated scalar, and the label beneath each space
specifies its vertex.

Using the weight-ordered bases
\[
\begin{aligned}
    \mathcal B_1&=(e_{1,0},e_{1,1}),\\
    \mathcal B_2&=(e_{2,0},e_{2,1}),\\
    \mathcal B_3&=(e_{3,0},e_{3,1},
        \varepsilon_3e_{3,0},\varepsilon_3e_{3,1}),\\
    \mathcal B_4&=(e_4,\varepsilon_4e_4,
        \varepsilon_4^2e_4,\varepsilon_4^3e_4),
\end{aligned}
\]
where $e_{i,a}=1\otimes1$ and $e_4=e_{4,0}$,
we obtain the following representation of $\Pi(C,D)$:
\[
\begin{tikzcd}[column sep=3em]
    \underset{1}{\mathbb C^2}
        \arrow[r,shift right=.6ex,"0"']
    & \underset{2}{\mathbb C^2}
        \arrow[l,shift right=.6ex,"X_{12}"']
        \arrow[r,shift left=.6ex,"X_{32}"]
    & \underset{3}{\mathbb C^4}
        \arrow[l,shift left=.6ex,"X_{23}"]
        \arrow[r,shift left=.6ex,"X_{43}"]
        \arrow[loop above,"N_3"]
    & \underset{4}{\mathbb C^4}
        \arrow[l,shift left=.6ex,"X_{34}"]
        \arrow[loop above,"N_4"]
\end{tikzcd}
\]
Here the loop operators are
\[
    N_3=
    \begin{pmatrix}
        0&0&0&0\\
        0&0&0&0\\
        1&0&0&0\\
        0&1&0&0
    \end{pmatrix},
    \qquad
    N_4=
    \begin{pmatrix}
        0&0&0&0\\
        1&0&0&0\\
        0&1&0&0\\
        0&0&1&0
    \end{pmatrix}.
\]
The loops at vertices $1$ and $2$ act as zero and are omitted from the diagram. With respect to the chosen bases and using column vectors, the arrow matrices are
\[
\begin{aligned}
    X_{12}
    &=
    \begin{pmatrix}
        x_0&0\\
        0&x_1
    \end{pmatrix},
    &
    X_{21}
    &=
    \begin{pmatrix}
        0&0\\
        0&0
    \end{pmatrix},
    \\[3mm]
    X_{32}
    &=
    \begin{pmatrix}
        0&0\\
        0&0\\
        y_0&0\\
        0&y_1
    \end{pmatrix},
    &
    X_{23}
    &=
    \begin{pmatrix}
        z_0&0&0&0\\
        0&z_1&0&0
    \end{pmatrix},
    \\[3mm]
    X_{34}
    &=
    \begin{pmatrix}
        0&0&0&0\\
        0&0&0&0\\
        u_0&0&0&0\\
        0&u_1&0&0
    \end{pmatrix},
    &
    X_{43}
    &=
    \begin{pmatrix}
        0&0&0&0\\
        0&0&0&0\\
        v_0&0&0&0\\
        0&v_1&0&0
    \end{pmatrix}.
\end{aligned}
\]
Here $X_{ij}$ represents the action of the arrow $\alpha_{ij}\colon j\to i$. The module $F(V)$ has rank vector $(2,2,2,1)$ and complex dimension vector $(2,2,4,4)$.
\end{example}

\subsection{Torus action and flag comparison}
\label{subsec:flag-comparison}
Let $T=\mathbb C^\times$ act on $F(V)$ through the weight decomposition by
\[
    g_i(t)v=t^w v
    \qquad
    (v\in F_i[w],\ t\in T).
\]
Write $g(t)=\bigoplus_{i\in I}g_i(t)$. 

A $\mathbb C$-linear map $\rho\colon F(V)_j\to F(V)_i$ is called homogeneous of degree $\delta\in\mathbb Z$ if
\[
    \rho(F_j[w])\subseteq F_i[w+\delta]
    \qquad\text{for every }w\in\mathbb Z.
\]

\begin{lemma}\label{lem:torus-action}
For every sequence $\mathbf i$, the weight action
on $F(V)$ induces an algebraic action of
$T=\mathbb C^\times$ on
$\mathrm{Fl}_{\mathbf i}^{E}(F(V))$, given by
\[
    t\cdot(U^\bullet)=g(t)U^\bullet.
\]
\end{lemma}

\begin{proof}
    The loop $\varepsilon_i$ has degree $m_i$, and the arrow degrees are
\[
    (\deg a_\ell,\deg b_\ell)=
    \begin{cases}
        (0,0),
            & d_i=d_j=1,\\
        (L/2,0),
            & d_i>1,\ d_j=1,\\
        (L-L/k_{ij},L-L/k_{ij}),
            & d_i,d_j>1.
    \end{cases}
\]
For every loop or arrow operator $\rho\colon F(V)_j\to F(V)_i$ of degree $\delta_\rho$, we have
\[
    g_i(t)\rho g_j(t)^{-1}
    =t^{\delta_\rho}\rho.
\]
Let $U=(U_i)_{i\in I}$ be a
$\Pi(C,D)$-submodule of $F(V)$. Then
\[
    \rho\bigl(g_j(t)U_j\bigr)
    =
    t^{-\delta_\rho}g_i(t)\rho(U_j)
    \subseteq g_i(t)U_i.
\]
Thus $g(t)U$ is again a submodule.

Let $U'\subseteq U$ be submodules with $U/U'\cong E_i$. The quotient $g(t)U/g(t)U'$ is supported at vertex $i$, so all arrows between distinct vertices act as zero. Under the linear identification induced by $g(t)$, its loop action is $t^{-m_i}$ times the original one, since
\[
    \varepsilon_i g(t)
    =t^{-m_i}g(t) \varepsilon_i.
\]
Rescaling the loop action by a nonzero scalar preserves the isomorphism class of $E_i$, so
\[
    g(t)U/g(t)U'\cong E_i.
\]
Applying this to each successive quotient shows that $g(t)$ preserves $\Fl_{\mathbf i}^{E}(F(V))$.
\end{proof}

Applying the Euler-characteristic fixed-point formula \cite[Corollary~2]{BB73} to the action in Lemma~\ref{lem:torus-action}, we obtain
\begin{equation}\label{eq:euler-fixed}
    \chi_c\bigl(\mathrm{Fl}_{\mathbf i}^{E}(F(V))\bigr)
    =
    \chi_c\bigl(\mathrm{Fl}_{\mathbf i}^{E}(F(V))^T\bigr).
\end{equation}
The fixed locus consists precisely of the flags whose terms are $T$-stable submodules. To describe these flags, we first identify the $T$-stable locally free submodules of $F(V)$.

\begin{lemma}\label{lem:free}
Let $U_i\subseteq F(V)_i$ be a $T$-stable $H_i$-submodule which is free over $H_i$. Then there are unique subspaces $W_{i,a}\subseteq V_{i,a}$ such that
\[
    U_i=
    \bigoplus_{a=0}^{m_i-1}
    H_i\otimes_{\mathbb C}W_{i,a}.
\]
Conversely, every submodule of this form is $T$-stable and free, and its quotient in $F(V)_i$ is free.
\end{lemma}

\begin{proof}
Since $U_i$ is $T$-stable, it decomposes into weight spaces:
\[
    U_i=\bigoplus_w\bigl(U_i\cap F_i[w]\bigr).
\]
The weights $a+pm_i$ are pairwise distinct, so there are unique subspaces
$W_{i,a,p}\subseteq V_{i,a}$ such that
\[
    U_i=
    \bigoplus_{a=0}^{m_i-1}
    \bigoplus_{p=0}^{d_i-1}
    \mathbb C\varepsilon_i^p
    \otimes_{\mathbb C}W_{i,a,p}.
\]
% Explicitly,
% \[
%     W_{i,a,p}
%     =
%     \left\{
%         v\in V_{i,a}
%         \ \middle|\
%         \varepsilon_i^p\otimes v\in U_i
%     \right\},
% \]
% so that
% \[
%     U_i\cap F_i[a+pm_i]
%     =
%     \mathbb C\varepsilon_i^p
%     \otimes_{\mathbb C}W_{i,a,p}.
% \]
Stability under $\varepsilon_i$ gives
\[
    W_{i,a,0}\subseteq W_{i,a,1}
    \subseteq\cdots\subseteq W_{i,a,d_i-1}
    \qquad(0\le a<m_i).
\]

Let $r=\rank_{H_i}U_i$. Since $U_i$ is free over $H_i$, we have
\[
    % \dim_{\mathbb C}U_i=d_i r,
    % \qquad
    r=\dim_{\mathbb C}
    \bigl(\varepsilon_i^{d_i-1}U_i\bigr)=\dim_{\mathbb C}\bigl( \bigoplus_{a=0}^{m_i-1}
    \mathbb C\varepsilon_i^{d_i-1}
    \otimes_{\mathbb C}W_{i,a,0} \bigr)=\sum_{a=0}^{m_i-1}
    \dim_{\mathbb C}W_{i,a,0}.
\]
% On the other hand, applying $\varepsilon_i^{d_i-1}$ to the weight decomposition above gives
% \[
%     \varepsilon_i^{d_i-1}U_i
%     =
%     \bigoplus_{a=0}^{m_i-1}
%     \mathbb C\varepsilon_i^{d_i-1}
%     \otimes_{\mathbb C}W_{i,a,0}.
% \]
% Thus
% \[
%     \sum_{a=0}^{m_i-1}
%     \dim_{\mathbb C}W_{i,a,0}=r.
% \]
Thus
\[
\begin{aligned}
    d_i r
    &=\dim_{\mathbb C}U_i\\
    &=\sum_{a=0}^{m_i-1}
      \sum_{p=0}^{d_i-1}
      \dim_{\mathbb C}W_{i,a,p}\\
    &\ge
      d_i\sum_{a=0}^{m_i-1}
      \dim_{\mathbb C}W_{i,a,0}
      =d_i r.
\end{aligned}
\]
Hence we have $W_{i,a,p}=W_{i,a,0}$ for all $a,p$. Taking $W_{i,a}=W_{i,a,0}$ proves the decomposition and uniqueness.
% Setting $W_{i,a}=W_{i,a,0}$, we obtain
% \[
%     U_i=
%     \bigoplus_{a=0}^{m_i-1}
%     H_i\otimes_{\mathbb C}W_{i,a}.
% \]
% These subspaces are uniquely determined by
% \[
%     U_i\cap F_i[a]
%     =
%     \mathbb C1\otimes_{\mathbb C}W_{i,a},
%     \qquad 0\le a<m_i.
% \]

Conversely, every submodule of the stated form is free over $H_i$ and $T$-stable. Moreover,
\[
    F(V)_i/U_i
    \cong
    \bigoplus_{a=0}^{m_i-1}
    H_i\otimes_{\mathbb C}(V_{i,a}/W_{i,a}),
\]
so the quotient is also free over $H_i$.
\end{proof}

We now describe the $T$-fixed flags in terms of flags of $Q$-subrepresentations.

\begin{lemma}\label{lem:flag-comparison}
Let $V$ be a finite-dimensional representation of $Q$, and let $\mathbf i=(i_1,\ldots,i_r)$. There is a natural isomorphism of varieties
\[
    \mathrm{Fl}_{\mathbf i}^{E}(F(V))^T
    \cong
    \bigsqcup_{\substack{0\le a_k<m_{i_k}\\1\le k\le r}}
    \mathrm{Fl}^{Q}_{((i_1,a_1),\ldots,(i_r,a_r))}(V).
\]
Consequently,
\[
    \chi_c\bigl(\mathrm{Fl}_{\mathbf i}^{E}(F(V))\bigr)
    =
    \sum_{\substack{0\le a_k<m_{i_k}\\1\le k\le r}}
    \chi_c\bigl(
        \mathrm{Fl}^{Q}_{((i_1,a_1),\ldots,(i_r,a_r))}(V)
    \bigr).
\]
\end{lemma}

\begin{proof}
Let $U^\bullet=\{0=U^0\subset U^1\subset\cdots\subset U^r=F(V)\} \in \Fl_{\mathbf i}^{E}(F(V))^T$. Each $U_i^k$ is $T$-stable and free over $H_i$. Lemma~\ref{lem:free} therefore gives unique subspaces $W_{i,a}^k\subseteq V_{i,a}$ such that
\[
    U_i^k=
    \bigoplus_{a=0}^{m_i-1}
    H_i\otimes_{\mathbb C}W_{i,a}^k.
\]
The inclusions $U^{k-1}\subseteq U^k$ imply $W_{i,a}^{k-1}\subseteq W_{i,a}^k$. By the construction of $F(V)$, the submodule condition on $U^k$ implies that $W^k=(W_{i,a}^k)$ is a $Q$-subrepresentation of $V$. Thus $W^\bullet=\{ 0=W^0\subset W^1\subset\cdots\subset W^r=V \}$ is a flag of $Q$-subrepresentations. Moreover,
\begin{align*}
    U_i^k/U_i^{k-1}
    \cong
    \bigoplus_{a=0}^{m_i-1}
    H_i\otimes_{\mathbb C}
    (W_{i,a}^k/W_{i,a}^{k-1}).
\end{align*}
Since $U^k/U^{k-1}\cong E_{i_k}$, there is a unique $0\le a_k<m_{i_k}$ for which $W_{i_k,a_k}^k/W_{i_k,a_k}^{k-1}$ is one-dimensional, and all other quotient spaces are zero. As $Q$ has no loops, we obtain
\begin{align*}
    W^k/W^{k-1}\cong S_{(i_k,a_k)}.
\end{align*}
Conversely, let $W^\bullet \in \Fl^{Q}_{((i_1,a_1),\ldots,(i_r,a_r))}(V)$. Let $U_i^k= \bigoplus_{a=0}^{m_i-1}H_i\otimes_{\mathbb C}W_{i,a}^k$.
By Lemma~\ref{lem:free}, each $U_i^k$ is a $T$-stable free $H_i$-submodule of $F(V)_i$. Since $W^k$ is a $Q$-subrepresentation, the construction of $F(V)$ shows that $U^k=(U_i^k)_{i\in I}$ is a $\Pi(C,D)$-submodule. The inclusions in $W^\bullet$ give
\begin{align*}
    0=U^0\subset U^1\subset\cdots\subset U^r=F(V).
\end{align*}
Since $W^k/W^{k-1}\cong S_{(i_k,a_k)}$, as $H_i$-modules we have
\[
\begin{aligned}
    U_i^k/U_i^{k-1}
    &\cong
    \bigoplus_{a=0}^{m_i-1}
    H_i\otimes_{\mathbb C}
    (W_{i,a}^k/W_{i,a}^{k-1})\\
    &\cong
    \begin{cases}
        H_i,&i=i_k,\\
        0,&i\ne i_k.
    \end{cases}
\end{aligned}
\]
Thus $U^k/U^{k-1}\cong E_{i_k}$, so $U^\bullet$ is a $T$-fixed $E$-flag. The two constructions are inverse by the uniqueness in Lemma~\ref{lem:free}.

It remains to show that both maps are regular.  For a fixed sequence $(a_1,\ldots,a_r)$, the dimensions of all spaces $W_{i,a}^k$ are fixed. The map $W^\bullet\mapsto U^\bullet$ is regular, since it is given by tensor products with the fixed spaces $H_i$ and direct sums.
Its inverse recovers $W_{i,a}^k$ from
\[
    U_i^k\cap F_i[a]
    =
    \mathbb C1\otimes_{\mathbb C}W_{i,a}^k.
\]
Since these intersections have fixed dimensions, the inverse map is also regular.

The dimensions $\dim_{\mathbb C}(U_i^k\cap F_i[w])$
are locally constant on the fixed locus of the
ambient flag variety.
Hence the loci indexed by $(a_1,\ldots,a_r)$
are open and closed in
$\mathrm{Fl}_{\mathbf i}^{E}(F(V))^T$. This proves the isomorphism of varieties.

Finally, the Euler-characteristic identity follows from \eqref{eq:euler-fixed} and additivity of $\chi_c$.
\end{proof}

We can now complete the proof of the main theorem.

\begin{proof}[Proof of Theorem~\ref{thm:main}]
By Corollary~\ref{cor:canonical-surjection}, $\Phi_{C,D}$ is a surjective rank-graded Hopf algebra homomorphism. Let $J:U(\mathfrak n_+(A))\hookrightarrow \mathcal M(Q)$ be the injection in Theorem~\ref{thm:auxiliary-embedding}. For $\mathbf i=(i_1,\ldots,i_r)$ and every $Q$-representation $V$,  Lemma~\ref{lem:flag-comparison} gives
\[
\begin{aligned}
    &(\theta_{i_1}*\cdots*\theta_{i_r})(F(V))\\
    &\qquad=
    \chi_c\bigl(\mathrm{Fl}_{\mathbf i}^{E}(F(V))\bigr)\\
    &\qquad=
    \sum_{\substack{0\le a_k<m_{i_k}\\1\le k\le r}}
    \chi_c\bigl(
        \mathrm{Fl}^{Q}_{((i_1,a_1),\ldots,(i_r,a_r))}(V)
    \bigr)\\
    &\qquad=
    \sum_{\substack{0\le a_k<m_{i_k}\\1\le k\le r}}
    (s_{(i_1,a_1)}*\cdots*s_{(i_r,a_r)})(V)\\
    &\qquad=
    (\sigma_{i_1}*\cdots*\sigma_{i_r})(V).
\end{aligned}
\]
By linearity, we obtain
\begin{equation}\label{eq:evaluation-comparison}
    \Phi_{C,D}(u)(F(V))=J(u)(V)
\end{equation}
for every $u\in U(\mathfrak n_+(A))$ and every finite-dimensional representation $V$ of $Q$.

Now let $u\in\ker\Phi_{C,D}$. Equation~\eqref{eq:evaluation-comparison} gives
\[
    J(u)(V)=\Phi_{C,D}(u)(F(V))=0
\]
for every finite-dimensional representation $V$ of $Q$. Thus $J(u)=0$, and injectivity of $J$ gives $u=0$. This proves that $\Phi_{C,D}$ is an isomorphism. Its restriction identifies $\nplus(A)$ with $\mathfrak l(C,D)$ since it sends $e_i$ to $\theta_i$.
\end{proof}



























%%%%%%%%%%%%%%%%%%%%%%%%%%%%%%%%%%%%%%%%%%%%%%%%%%%%%%%%%%%%%%%%%%%%%%%%%
\section{Nonrealization and Serre quotients}\label{app:serre-quotients}
Put $A=A(C,D)$ and let
\begin{equation}\label{eq:app-exceptional-edges}
    \mathscr B=
    \bigl\{
        \{i,j\}:c_{ij}<0,\ d_i,d_j>1,\
        \gcd(d_i,d_j)=1
    \bigr\}.
\end{equation}
Thus condition~\eqref{eq:edge-condition} holds if and only if $\mathscr B=\varnothing$. To prove Theorem~\ref{thm:nonrealization}, we first show that
$(\ad\theta_i)^{-a_{ij}}\theta_j\ne0$ for all $i\ne j$. We then construct a nonzero element of $\ker\Phi_{C,D}$ for each edge in $\mathscr B$.

\subsection{Proof of Theorem~\ref{thm:nonrealization}}
\begin{proposition}\label{prop:sharp-Serre}
For all distinct vertices $i,j\in I$, we have
\[
    (\ad\theta_i)^{-a_{ij}}(\theta_j)\ne0.
\]
\end{proposition}
\begin{proof}
    If $c_{ij}=0$, the assertion is immediate. If the restriction of $(C,D)$ to $\{i,j\}$ satisfies condition~\eqref{eq:edge-condition}, the assertion follows from Theorem~\ref{thm:main}. Therefore, we assume that 
    \[
    c_{ij}<0,\qquad p=d_i>1,\qquad q=d_j>1,\qquad \gcd(p,q)=1.
    \]
Put $g=g_{ij}$ and $N=|c_{ij}|=gq$. Following the construction in \cite[Lemma~6.1]{GLS18},
we construct a locally free $\Pi$-module $X$ on which $(\ad\theta_i)^{-a_{ij}}\theta_j$ is nonzero.

Take
\[
    X_i=P\oplus Q,\qquad P\cong Q\cong H_i^N,
    \qquad X_j=H_j,
\]
with $H_i$-bases $u_{\ell,t}$ of $P$ and $v_{\ell,t}$ of $Q$, where $1\le\ell\le g$ and $0\le t<q$.

Write
\[
    u_{\ell,t}^{(s)}=\varepsilon_i^s u_{\ell,t},
    \qquad
    v_{\ell,t}^{(s)}=\varepsilon_i^s v_{\ell,t}
    \qquad(0\le s<p),
\]
for the corresponding $\C$-bases of $P$ and $Q$, and let 
\[
    w_t=\varepsilon_j^t \qquad (0\le t<q)
\]
be the $\C$-basis of $X_j=H_j$.

The loops act by
\[
    \varepsilon_i u_{\ell,t}^{(s)}
    =
    \begin{cases}
        u_{\ell,t}^{(s+1)}, & 0\le s<p-1,\\
        0, & s=p-1,
    \end{cases}
    \qquad
    \varepsilon_i v_{\ell,t}^{(s)}
    =
    \begin{cases}
        v_{\ell,t}^{(s+1)}, & 0\le s<p-1,\\
        0, & s=p-1,
    \end{cases}
\]
and
\[
    \varepsilon_j w_t
    =
    \begin{cases}
        w_{t+1}, & 0\le t<q-1,\\
        0, & t=q-1.
    \end{cases}
\]
Define the arrow maps by
\[
    a_\ell(w_t)=u_{\ell,t}^{(p-1)},
    \qquad
    b_\ell(v_{\ell,t}^{(0)})=w_t
    \qquad(1\le\ell\le g,\ 0\le t<q),
\]
and all other basis vectors are mapped to zero. It follows that
\[
    \varepsilon_i^p=\varepsilon_j^q=0,\qquad
    b_\ell a_\ell=0,\qquad
    \varepsilon_i a_\ell=b_\ell\varepsilon_i=0.
\]
Thus these maps define a locally free $\Pi$-module $X$.

Since $\Ext^1_\Pi(E_i,E_i)=0$ and $\chi_c\bigl(\Fl^E_{(i,\ldots,i)}(E_i^{\oplus r})\bigr)=r!$, we have $\theta_i^{*r}=r!\,\mathbf1_{E_i^{\oplus r}}$. We now compute
\begin{align*}
    \bigl((\ad\theta_i)^{2N}\theta_j\bigr)(X)
    &=
    \sum_{s=0}^{2N}(-1)^{2N-s}\binom{2N}{s}
    \bigl(\theta_i^{*s}*\theta_j*
          \theta_i^{*(2N-s)}\bigr)(X)\\
    &=
    (2N)!\sum_{s=0}^{2N}(-1)^s
    \bigl(
        \mathbf1_{E_i^{\oplus s}}
        *\mathbf1_{E_j}
        *\mathbf1_{E_i^{\oplus(2N-s)}}
    \bigr)(X).
\end{align*}
For each $s$, we consider partial flags
\[
    0\subseteq(U,0)\subseteq(U,H_j)\subseteq X,
\]
where $U$ is a free $H_i$-direct summand of $X_i$ of rank $s$. By the definitions of $a_\ell$ and $b_\ell$, $(U,0)$ and $(U,H_j)$ are submodules of $X$ if and only if
\[
    \varepsilon_i^{p-1}P
        \subseteq U\subseteq P\oplus\varepsilon_iQ.
\]
Since $U$ is a direct summand of $X_i$, the first inclusion gives
\[
    N=\dim_\C\varepsilon_i^{p-1}P
      \le\dim_\C(U \cap \varepsilon_i^{p-1}X_i)
      =\dim_\C\varepsilon_i^{p-1}U=s.
\]
The second inclusion gives
\[
    s=\dim_\C(U/\varepsilon_iU)\le \dim_\C(P/\varepsilon_iP)=N.
\]
Hence $s=N$. 

Write $\pi_P$ and $\pi_Q$ for the projections from $P\oplus Q$ onto $P$ and $Q$, respectively. The projection $\pi_P|_U:U\to P$ induces an isomorphism
\[
    U/\varepsilon_iU \xrightarrow{\sim} P/\varepsilon_iP.
\]
By Nakayama's lemma, $\pi_P|_U$ is surjective, so it is an isomorphism.

Let $\mathcal U$ be the variety of these submodules. For $\varphi\in\Hom_{H_i}(P,\varepsilon_iQ)$, set
\[
    U_\varphi=\{x+\varphi(x)\mid x\in P\}.
\]
The map $x\mapsto x+\varphi(x)$ is an $H_i$-linear isomorphism, so $U_\varphi$ is free of rank $N$. It follows that
\[
    \varepsilon_i^{p-1}P \subseteq U_\varphi \subseteq P\oplus\varepsilon_iQ.
\]
Therefore the map
\[
    F:\Hom_{H_i}(P,\varepsilon_iQ)
    \longrightarrow\mathcal U,
    \qquad \varphi\longmapsto U_\varphi,
\]
is well-defined.

Conversely, define
\[
    G: \mathcal U
    \longrightarrow \Hom_{H_i}(P,\varepsilon_iQ),
    \qquad U\longmapsto\pi_Q\circ(\pi_P|_U)^{-1}.
\]
The maps $F$ and $G$ are mutually inverse and regular in the Grassmannian chart of complements to $Q$. Thus
\[
    \mathcal U\cong\Hom_{H_i}(P,\varepsilon_iQ)
        \cong\mathbb A^{(p-1)N^2}.
\]
Since $\chi_c(\mathbb A^{(p-1)N^2})=1$, we obtain
\[
    \bigl((\ad\theta_i)^{-a_{ij}}\theta_j\bigr)(X)
      =(-1)^N(2N)!\ne0.
\]
\end{proof}

We now show that each edge in $\mathscr B$ gives a nonzero element of $\ker\Phi_{C,D}$.

\begin{proposition}\label{prop:app-extra-relations}
Let $\{i,j\}\in\mathscr B$, and set
$a=-a_{ij}$ and $b=-a_{ji}$. Then
\[
    (\ad\theta_i)^{ab-2}(\ad\theta_j)^b\theta_i=0,
\]
whereas
\[
    (\ad e_i)^{ab-2}(\ad e_j)^b e_i\ne0
    \qquad\text{in }\mathfrak n_+(A).
\]
In particular, if condition~\eqref{eq:edge-condition}
fails, then $\Phi_{C,D}$ is not injective.
\end{proposition}

\begin{proof}
    Since $\gcd(d_i,d_j)=1$, we have $a=2g_{ij}d_j$ and $b=2g_{ij}d_i$. Put $N=b|c_{ij}|=ab/2$, and set
    \[
    u=(\ad\theta_j)^b \theta_i, \qquad z=(\ad \theta_i)^{ab-2}u.
    \]
    Since $z$ is primitive and homogeneous of rank $(ab-1)\mathbf e_i+b\mathbf e_j$, it suffices to show that $z(X)=0$ for every indecomposable locally free module $X$ of this rank. Then we have
    \[
    z=\sum_{k=0}^{ab-2}(-1)^k\binom{ab-2}{k}\theta_i^{*(ab-2-k)}*u*\theta_i^{*k}.
    \]
    For each $k$, consider a flag
    \[
    0\subseteq X'\subseteq X''\subseteq X
    \]
such that $X'\cong E_i^{\oplus(ab-2-k)}$, $X/X''\cong E_i^{\oplus k}$ and $Y=X''/X'$ is locally free of rank $\mathbf e_i+b\mathbf e_j$. We show that $E_j$ is a direct summand of every such $Y$. This will give $u(Y)=0$ by Lemma~\ref{lem:primitive-support} and hence $z(X)=0$.

Write $\alpha=X_{i,\mathrm{in}}$ and $\beta=X_{i,\mathrm{out}}$. Then
     \[
     \widetilde X_i\cong H_i^N,\qquad X_i\cong H_i^{2N-1},
    \qquad
    \alpha:H_i^N\longrightarrow H_i^{2N-1},\quad
    \beta:H_i^{2N-1}\longrightarrow H_i^N,\quad
    \alpha\beta=0.
     \]

By Smith normal form, each of $\Im \alpha$ and $\Im \beta$ can be generated by at most $N$ elements. If both could be generated by at most $N-1$ elements, then Lemma~\ref{lem:general splitting} would give $E_i$ as a direct summand of $X$, since $2N-1>2(N-1)$. This contradicts the indecomposability of $X$.

Thus at least one of $\Im \alpha$ and $\Im \beta$ is minimally generated by $N$ elements. Suppose first that $\Im \alpha$ is minimally generated by $N$ elements as an $H_i$-module.

Let $u_1,\ldots,u_N$ be an $H_i$-basis of $\widetilde X_i$. Then $\Im \alpha=\sum_{r=1}^N H_i \alpha(u_r)$. We claim that $\Ker \alpha \subseteq \varepsilon_i \widetilde X_i$. Indeed, for $v=\sum_{r=1}^N c_r u_r \in \Ker \alpha$, we have
\[
    \sum_{r=1}^N c_r\alpha(u_r)=0.
\]
If some $c_r$ were invertible, then $\alpha(u_r)$ would be generated by the remaining $N-1$ elements, contradicting the minimality of $N$. 
Thus every $c_r$ belongs to $\varepsilon_iH_i$, and hence $v\in\varepsilon_i\widetilde X_i$, as claimed. 
Since $\alpha\beta=0$, we obtain
\[
    \Im\beta\subseteq\Ker\alpha
    \subseteq\varepsilon_i\widetilde X_i.
\]
Consequently, 
\[
    \beta(\varepsilon_i^{d_i-1}X_i)
    =\varepsilon_i^{d_i-1}\beta(X_i)=0.
\]
Similarly, suppose that $\Im\beta$ is minimally generated by $N$ elements as an $H_i$-module. By the same proof we have $\Im\alpha\subseteq\varepsilon_i X_i$.

Thus $X$ satisfies at least one of
\[
    \text{(a)}\quad
    \Im X_{i,\mathrm{in}} \subseteq\varepsilon_iX_i,
    \qquad
    \text{(b)}\quad X_{i,\mathrm{out}} (\varepsilon_i^{d_i-1}X_i)=0.
\]
Now consider the flag as above and put $Y=X''/X'$. We have $X'_j=0$ and $X''_j=X_j$.

In case (a), since $X''$ is a submodule and $X''_j=X_j$, we have $\Im\alpha\subseteq X''_i$. Moreover, $X_i/X''_i$ is free over $H_i$, so
\[
    X''_i\cap\varepsilon_iX_i=\varepsilon_iX''_i.
\]
Thus $\Im\alpha\subseteq \varepsilon_iX''_i$. Passing to the quotient by $X'_i$ gives 
\[
    \Im Y_{i,\mathrm{in}}\subseteq\varepsilon_iY_i.
\]
In case (b), since $X'$ is a submodule and $X'_j=0$, we have $\beta(X'_i)=0$. Thus $\beta|_{X''_i}$ induces $Y_{i,\mathrm{out}}$ on $Y_i=X''_i/X'_i$. Since $\beta(\varepsilon_i^{d_i-1}X_i)=0$, we obtain
\[
    Y_{i,\mathrm{out}}(\varepsilon_i^{d_i-1}Y_i)=0.
\]
We show that either condition makes $Y$ decomposable. Since $Y$ is locally free of rank $\mathbf e_i+b\mathbf e_j$, we have
\[
    Y_i\cong H_i,\qquad Y_j\cong H_j^b,
    \qquad
    \widetilde Y_j
    ={}_jH_i\otimes_{H_i}Y_i
    \cong H_j^{|c_{ji}|}.
\]
Thus these maps have the form
\[
    \widetilde Y_j
    \xrightarrow{\,Y_{j,\mathrm{in}}\,}
    Y_j
    \xrightarrow{\,Y_{j,\mathrm{out}}\,}
    \widetilde Y_j.
\]
Since $f_{ji}=d_i$ and $f_{ij}=d_j$, all commutation relations $\varepsilon_j^{f_{ij}}\alpha_{ji}^{(\ell)}=\alpha_{ji}^{(\ell)}\varepsilon_i^{f_{ji}}$ hold with both sides equal to zero. Hence
\[
    {}_jH_i\cong
    (H_j\otimes_{\mathbb C}H_i)^{\oplus g_{ij}}
\]
as $(H_j,H_i)$-bimodules, and therefore
\[
    \widetilde Y_j
    ={}_jH_i\otimes_{H_i}Y_i
    \cong H_j\otimes_{\mathbb C}Y_i^{\oplus g_{ij}}.
\]

In case (a), each arrow map from $j$ to $i$ sends $Y_j$ into $\varepsilon_iY_i$. The definition of $Y_{j,\mathrm{out}}$ gives
\[
    \Im Y_{j,\mathrm{out}}
    \subseteq H_j\otimes_{\mathbb C}
    (\varepsilon_iY_i)^{\oplus g_{ij}}.
\]
The module on the right is free over $H_j$ of rank $g_{ij}(d_i-1)=|c_{ji}|-g_{ij}$. Hence $\Im Y_{j,\mathrm{out}}$ can be generated by at most $|c_{ji}|-g_{ij}$ elements as an $H_j$-module.

In case (b), each arrow map from $i$ to $j$ vanishes on $\varepsilon_i^{d_i-1}Y_i$. The definition of $Y_{j,\mathrm{in}}$ gives
\[
    H_j\otimes_{\mathbb C}
    (\varepsilon_i^{d_i-1}Y_i)^{\oplus g_{ij}}
    \subseteq\Ker Y_{j,\mathrm{in}}.
\]
The module on the left is free over $H_j$ of rank $g_{ij}$. Hence $\Im Y_{j,\mathrm{in}}$ can be generated by at most $|c_{ji}|-g_{ij}$ elements as an $H_j$-module.

In either case, one of the two images can be generated by at most $|c_{ji}|-g_{ij}$ elements, and the other by at most $|c_{ji}|$ elements, as $H_j$-modules. Since
\[
    \rank_{H_j}Y_j=b=2|c_{ji}|
    >|c_{ji}|+(|c_{ji}|-g_{ij}),
\]
Lemma~\ref{lem:general splitting} gives $E_j$ as a direct summand of $Y$. Since $Y_i\ne0$, the module $Y$ is decomposable. Thus $u(Y)=0$ and it follows that $z=0$.

Finally, to prove that $(\ad e_i)^{ab-2}(\ad e_j)^b e_i\ne0$, we apply
$(\ad f_j)^b(\ad f_i)^{ab-2}$ to this element. Since $b\ge2$, we have
\begin{align*}
    [f_i,(\ad e_j)^b e_i]&=-(\ad e_j)^b h_i=0,\\
    [h_i,(\ad e_j)^b e_i] &=(2-ab)(\ad e_j)^b e_i.
\end{align*}
Hence we have
\begin{align*}
    (\ad f_j)^b(\ad f_i)^{ab-2}(\ad e_i)^{ab-2}(\ad e_j)^b e_i&=((ab-2)!)^2
  (\ad f_j)^b(\ad e_j)^b e_i\\
&=((ab-2)!)^2(b!)^2e_i\ne0.
\end{align*}
\end{proof}

\begin{proof}[Proof of Theorem~\ref{thm:nonrealization}]
Assume that $(C,D)$ does not satisfy condition~\eqref{eq:edge-condition}. Put $L=\mathfrak l(C,D)$. Suppose that $L\cong\mathfrak n_+(B)$ for some symmetrizable generalized Cartan matrix $B$. 

Let $\{L^n\}_{n\ge1}$ denote the lower central series of $L$, where $L^1=L$ and $L^n=[L,L^{n-1}]$ for $n\ge2$. We use the same notation for $\mathfrak n_+(A)$ and $\mathfrak n_+(B)$. By Proposition~\ref{prop:app-extra-relations}, the canonical surjection
\[
    \pi:\mathfrak n_+(A)\longrightarrow L,
    \qquad e_i\longmapsto\theta_i,
\]
is not injective. Choose a nonzero homogeneous element $z\in\ker\pi$, and let $d_z$ be its total degree. Let $N=\max\left\{ d_z,\ \max_{i\ne j}(1-a_{ij}),\ \max_{i\ne j}(1-b_{ij})\right\}$ and set 
\[
L_N=L/L^{N+1}, \qquad A_N=\nplus(A)/\nplus(A)^{N+1} ,\qquad B_N=\nplus(B)/\nplus(B)^{N+1}.
\]
These are finite-dimensional nilpotent Lie algebras, and the assumed isomorphism induces $L_N\cong B_N$.

Write $\bar\theta_i$ for the image of $\theta_i$ in $L_N$. The rank grading of $L$ descends to $L_N$ and defines commuting semisimple derivations $D_i$, $i\in I$, satisfying
\[
    D_i(\bar\theta_j)=\delta_{ij}\bar\theta_j.
\]
Thus
\[
    T=\bigoplus_{i\in I}\C D_i
\]
is a torus of dimension $|I|$. Since $\dim_\C L_N/[L_N,L_N]=|I|$, it follows from
\cite[Lemma~2.8]{San82} and \cite[Definition~2.9]{San82} that $T$ is maximal and $L_N$ is of maximal rank. $(\bar\theta_i)_{i\in I}$ is a $T$-msg in the sense of \cite[Definition~2.5]{San82}.

By Proposition~\ref{prop:A-Serre} and Proposition~\ref{prop:sharp-Serre}, we have
\[
    (\ad\bar\theta_i)^{-a_{ij}}\bar\theta_j\ne0,
    \qquad
    (\ad\bar\theta_i)^{1-a_{ij}}\bar\theta_j=0
    \qquad(i\ne j).
\]
Thus $A$ is the  Cartan matrix associated with $L_N$ in the sense of \cite[Theorem~3.2]{San82}. The same argument applied to $\nplus(B)$ shows that $B_N$ is of maximal rank with associated Cartan matrix $B$. Since $L_N\cong B_N$, \cite[Theorem~3.2]{San82} implies that $A$ and $B$ differ by a simultaneous permutation.

The map $\pi$ induces a surjection $\bar\pi: A_N \twoheadrightarrow L_N$. Let $\bar z$ be the image of $z$ in $A_N$. We have $\bar z\ne0$, while $\bar\pi(\bar z)=0$. Thus $\bar\pi$ is not injective. Therefore
\[
    \dim_\C L_N< \dim_\C A_N =\dim_\C B_N,
\]
contrary to $L_N\cong B_N$. This proves the assertion for $L$.

Lemma~\ref{lem:enveloping-realization} gives $\mathcal M(C,D)\cong U(L)$ as Hopf algebras. Suppose that $\mathcal M(C,D)\cong U(\mathfrak n_+(B))$ as unital $\C$-algebras for some symmetrizable generalized Cartan matrix $B$. Then $U(L)\cong U(\mathfrak n_+(B))$ as unital $\C$-algebras, and \cite[Proposition~4.1]{RU07} implies that
\[
    \operatorname{gr}L
    \cong\operatorname{gr}\mathfrak n_+(B),
\]
where $\operatorname{gr}L=\bigoplus_{r\ge1}L^r/L^{r+1}$. Since $L$ and $\nplus(B)$ are positively graded and generated in degree one, their total gradings give Lie algebra isomorphisms
\[
    L\cong\operatorname{gr}L,\qquad
    \mathfrak n_+(B)\cong\operatorname{gr}\mathfrak n_+(B).
\]
Hence $L\cong\mathfrak n_+(B)$, contradicting the assertion for $L$ proved above.
\end{proof}


\subsection{Serre quotient realizations}
We now return to an arbitrary pair $(C,D)$. Put $\mathcal M=\mathcal M(C,D)$ and let $\mathscr B$ be as in~\eqref{eq:app-exceptional-edges}.
Define $B=B(C,D)=(b_{ij})_{i,j\in I}$ by
\begin{equation}\label{eq:app-matrix-B}
 b_{ii}=2,\qquad
 b_{ij}=
 \begin{cases}
 c_{ij},&\{i,j\}\in\mathscr B,\\
 a_{ij},&\{i,j\}\notin\mathscr B.
 \end{cases}
\end{equation}
The positive integral diagonal matrix
\[
 D_A=\diag(2d_i/\lambda_i)_{i\in I}
\]
symmetrizes both $A$ and $B$. In particular, $B$ is a symmetrizable generalized Cartan matrix.

Let $\mathcal J_{\mathscr B}$ be the two-sided ideal of $\mathcal M$ generated by
\begin{equation}\label{eq:app-exceptional-ideal}
 (\ad\theta_i)^{1-c_{ij}}\theta_j
\end{equation}
for all ordered pairs $(i,j)$ with $\{i,j\}\in\mathscr B$. Since these elements are primitive, $\mathcal J_{\mathscr B}$ is a Hopf ideal.

\begin{lemma}[An intermediate Serre quotient] \label{lem:app-intermediate-quotient}
For every pair $(C,D)$, there is a canonical rank-graded Hopf algebra isomorphism
\begin{equation}\label{eq:app-intermediate-isomorphism}
\begin{aligned}
 U(\nplus(B))&\xrightarrow{\sim}
\mathcal M(C,D)/\mathcal J_{\mathscr B},\\
 e_i&\longmapsto\theta_i+\mathcal J_{\mathscr B}.
\end{aligned}
\end{equation}
\end{lemma}

\begin{proof}
    The $B$-Serre relations follow from Proposition~\ref{prop:A-Serre} and the definition of $\mathcal J_{\mathscr B}$. Thus the assignment in the statement induces a rank-graded Hopf algebra homomorphism
\[
    \psi_B:U(\nplus(B))
    \longrightarrow
    \mathcal M(C,D)/\mathcal J_{\mathscr B},
    \qquad
    e_i\longmapsto\theta_i+\mathcal J_{\mathscr B}.
\]
Since the classes $\theta_i+\mathcal J_{\mathscr B}$ generate the quotient as a unital algebra, $\psi_B$ is surjective.

For injectivity, we modify the auxiliary quiver $Q$ of Section~\ref{sec:injectivity}. Keep its vertex set and all arrows on edges outside $\mathscr B$. For each $\{i,j\}\in\mathscr B$, choose a direction $j\to i$ and retain only the arrows
\[
    \alpha_{ij,w}^{(\ell)}:
    (j,[w]_{m_j})\longrightarrow(i,[w]_{m_i}),
    \qquad
    0\le w<L,\quad 1\le\ell\le g_{ij}.
\]
Omit the arrows $\beta_{ij,w}^{(\ell)}$. Denote the resulting quiver by $Q'$ and its symmetric generalized Cartan matrix by $\widetilde B$.

For $\{i,j\}\in\mathscr B$, we have $f_{ij}=d_j$ and $f_{ji}=d_i$. For each fixed vertex $(j,b)$, it follows that
\[
    \sum_{a=0}^{m_i-1}
        \widetilde b_{(i,a),(j,b)}
    =-g_{ij}\frac{L}{m_j}
    =-g_{ij}d_j
    =c_{ij}
    =b_{ij}.
\]
Interchanging $i$ and $j$ gives the corresponding identity for $b_{ji}$. For all other edges, the arrows are unchanged, so the folding identities follow from Section~\ref{sec:injectivity}.

$Q'$ satisfies the hypotheses of Lemma~\ref{lem:inj-cover}. Combining the resulting embedding with Theorem~\ref{thm:inj-hall}, we obtain an injective algebra homomorphism
\[
\begin{aligned}
    J_B:U(\nplus(B))&\hookrightarrow\mathcal M(Q'),\\
    e_i&\longmapsto
    \sigma_i:=
    \sum_{a=0}^{m_i-1}\mathbf1_{S_{(i,a)}}.
\end{aligned}
\]
Regard a representation $V$ of $Q'$ as a representation of $Q$ by assigning zero maps to the deleted arrows. Define $F(V)$ using the formulas of Section~\ref{sec:injectivity}. For $\{i,j\}\in\mathscr B$, we have $k_{ij}=1$, so $t=L$, the weight shift is zero, and $b_\ell=0$. Since $f_{ji}=d_i$ and $f_{ij}=d_j$, the commutation relations on this edge hold by loop nilpotence, and its mesh contributions vanish. For all other edges, the verification in Lemma~\ref{lem:relations} applies unchanged. Thus $F(V)$ is a locally free $\Pi(C,D)$-module, and all its operators remain homogeneous.

By the same argument as in the proof of Lemma~\ref{lem:flag-comparison}
(see also \cite[Appendix~A.2]{Fei26}), we obtain
\[
    (\theta_{i_1}*\cdots*\theta_{i_r})(F(V))
    =
    (\sigma_{i_1}*\cdots*\sigma_{i_r})(V)
\]
for every word and every representation $V$ of $Q'$. Consequently, the assignment $\theta_i\mapsto\sigma_i$ induces an algebra homomorphism
\[
    \rho:\mathcal M(C,D)\longrightarrow\mathcal M(Q').
\]
Since the $\sigma_i$ satisfy the $B$-Serre relations, $\rho$ annihilates $\mathcal J_{\mathscr B}$ and induces a homomorphism
\[
    \bar\rho:
    \mathcal M(C,D)/\mathcal J_{\mathscr B}
    \longrightarrow\mathcal M(Q').
\]
By construction, $\bar\rho\circ\psi_B=J_B$. Since $J_B$ is injective, $\psi_B$ is injective. This proves the claimed isomorphism.
\end{proof}

As a consequence, we recover the full $C$-Serre quotient realization established by Fei \cite[Theorem~10.1]{Fei26}. In particular, this gives the algebra isomorphism in \cite[Theorem~7.4(i)]{GLS18} without assuming the Support Conjecture.

\begin{corollary}[The full $C$-Serre quotient] \label{cor:app-full-Serre-quotient}
Let $\mathcal I_C$ be the two-sided ideal of $\mathcal M(C,D)$ generated by
\[
    (\ad\theta_i)^{1-c_{ij}}\theta_j,
    \qquad i\ne j.
\]
For every pair $(C,D)$, there is a canonical rank-graded Hopf algebra isomorphism
\[
\begin{aligned}
    U(\nplus(C))
    &\xrightarrow{\sim}
    \mathcal M(C,D)/\mathcal I_C,\\
    e_i&\longmapsto\theta_i+\mathcal I_C.
\end{aligned}
\]
\end{corollary}

\begin{proof}
Since $\mathcal J_{\mathscr B}\subseteq\mathcal I_C$ and $b_{ij}\le c_{ij}$ for $i\ne j$, the assertion follows from Lemma~\ref{lem:app-intermediate-quotient} and the Serre presentation of $U(\nplus(C))$.
\end{proof}




%%%%%%%%%%%%%%%%%%%%%%%%%%%%%%%%%%%%%%%%%%
\section*{Declaration of AI-assisted research and manuscript preparation}
The author used ChatGPT for literature searches, language editing, computations in examples, and discussions of proof strategies. This assistance helped identify a counterexample to an earlier formulation of Theorem~\ref{thm:main}. Further computations led to Proposition~\ref{prop:app-extra-relations}, which the author used to prove Theorem~\ref{thm:nonrealization}. The author has carefully checked the paper and takes full responsibility for its correctness. 
%%%%%%%%%%%%%%%%%%%%%%%%%%%%%%%%%%%%%%%%%%
\begin{thebibliography}{SHI26}

    \bibitem[BB73]{BB73} A.~Bia{\l}ynicki-Birula, \emph{On fixed point schemes of actions of multiplicative and additive groups}, Topology \textbf{12} (1973), 99--103.

    \bibitem[CG25]{CG25} A.~Castillo G\'omez and C.~Geiss, \emph{A geometric construction of $U(\mathfrak n)$ for affine Kac--Moody algebras of type $\widetilde{\mathsf C}_n$}, J. Pure Appl. Algebra \textbf{229} (2025), no.~10, Paper No.~108069.

    \bibitem[Dem11]{Dem11} L.~Demonet, \emph{Categorification of skew-symmetrizable cluster algebras}, Algebr. Represent. Theory \textbf{14} (2011), no.~6, 1087--1162.

    \bibitem[Fei26]{Fei26} J.~Fei, \emph{Ramified convolution and monodromy bases for symmetrizable Kac--Moody algebras}, \arxiv{2610.29051}, 2026.

    %\bibitem[GLS05]{GLS05} C.~Geiss, B.~Leclerc, and J.~Schr\"oer, \emph{Semicanonical bases and preprojective algebras}, Ann. Sci. \'Ecole Norm. Sup. (4) \textbf{38} (2005), no.~2, 193-253.

    %\bibitem[GLS06]{GLS06} C.~Geiss, B.~Leclerc, and J.~Schr\"oer, \emph{Rigid modules over preprojective algebras}, Invent. Math. \textbf{165} (2006), no.~3, 589-632.

    %\bibitem[GLS07]{GLS07} C.~Geiss, B.~Leclerc, and J.~Schr\"oer, \emph{Semicanonical bases and preprojective algebras II: A multiplication formula}, Compos. Math. \textbf{143} (2007), no.~5, 1313-1334.

    %\bibitem[GLS08]{GLS08} C. Gei\ss, B. Leclerc and J. Schr\"{o}er,c\emph{Partial flag varieties and preprojective algebras}, Ann. Inst. Fourier (Grenoble) \textbf{58} (2008), 825-876.

    %\bibitem[GLS11]{GLS11} C.~Gei\ss, B.~Leclerc and J.~Schr\"oer, \emph{Kac-Moody groups and cluster algebras}, Adv. Math. \textbf{228} (2011), no.~1, 329-433.

    \bibitem[GLS16]{GLS16} C. Gei\ss, B. Leclerc and J. Schr\"{o}er, \emph{Quivers with relations for symmetrizable Cartan matrices III: Convolution algebras}, Represent. Theory \textbf{20} (2016), 375--413. \arxiv{1511.06216}.

    \bibitem[GLS17]{GLS17} C. Gei\ss, B. Leclerc and J. Schr\"{o}er,
    \emph{Quivers with relations for symmetrizable Cartan matrices I:
    Foundations}, Invent. Math. \textbf{209} (2017), 61--158.

    \bibitem[GLS18]{GLS18} C. Gei\ss, B. Leclerc and J. Schr\"{o}er,
    \emph{Quivers with relations for symmetrizable Cartan matrices IV:
    Crystal graphs and semicanonical functions}, Selecta Math. (N.S.)
    \textbf{24} (2018), 3283--3348.

    \bibitem[Gr95]{Gr95} J.~A.~Green, \emph{Hall algebras, hereditary algebras and quantum groups}, Invent. Math. \textbf{120} (1995), no.~2, 361--377.

    \bibitem[Hub02]{Hub02} A.~W.~Hubery,
    \emph{Representations of quivers respecting a quiver automorphism and a theorem of Kac}, Ph.D. thesis, University of Leeds, 2002.

    \bibitem[JS12]{JS12} Y.~Jiang and J.~Sheng,
    \emph{Tame quivers and affine enveloping algebras}, Pure Appl. Math. Q. \textbf{8} (2012), no.~3, 609--651.

    \bibitem[Lus91]{Lus91} G.~Lusztig, \emph{Quivers, perverse sheaves, and quantized enveloping algebras}, J. Amer. Math. Soc. \textbf{4} (1991), no.~2, 365--421.

    \bibitem[Lus00]{Lus00} G.~Lusztig,
    \emph{Semicanonical bases arising from enveloping algebras}, Adv. Math. \textbf{151} (2000), no.~2, 129--139.

    \bibitem[Rin90]{Rin90} C.~M.~Ringel, \emph{Hall algebras and quantum groups}, Invent. Math. \textbf{101} (1990), 583--591.

    \bibitem[RU07]{RU07} D.~M.~Riley and H.~Usefi, \emph{The isomorphism problem for universal enveloping algebras of Lie algebras}, Algebr. Represent. Theory \textbf{10} (2007), no.~6, 517--532.

    \bibitem[San82]{San82} L.~J.~Santharoubane, \emph{Kac--Moody Lie algebras and the classification of nilpotent Lie algebras of maximal rank}, Canad. J. Math. \textbf{34} (1982), no.~6, 1215--1239.

    \bibitem[Sch91]{Sch91} A.~Schofield, \emph{Notes on constructing Lie algebras from finite-dimensional algebras}, preprint, 1991.

\end{thebibliography}
\end{document}

